% concatenated from Paper.tex
\documentclass[12pt, a4paper]{article}
\usepackage[left=2cm, right=2cm, top=2.5cm, bottom=2.5cm, headheight=15pt, includehead, includefoot]{geometry}
\usepackage{parskip} 
\setlength{\parindent}{15pt}

%% Required Packages
\usepackage{amsmath, amsthm, amssymb, bm}
\usepackage{mathrsfs}
\usepackage[noend]{algpseudocode}
\usepackage{graphicx, subcaption}
\usepackage{comment}
\usepackage{xcolor}
\usepackage[colorlinks=true, citecolor=blue, linkcolor=blue]{hyperref}
\usepackage{url} 
\usepackage{footmisc}
\usepackage{natbib}
\usepackage{authblk}
\usepackage[ruled]{algorithm2e}
\usepackage{colortbl}
\definecolor{lightteal}{RGB}{224,245,245}
\usepackage{enumitem}

%% Theorem Styles
\theoremstyle{plain}
\newtheorem{theorem}{Theorem}
\newtheorem{lemma}{Lemma}
\newtheorem{proposition}{Proposition}
\newtheorem{corollary}{Corollary}

\theoremstyle{definition}
\newtheorem{definition}{Definition}
\newtheorem{example}{Example}

\theoremstyle{remark}
\newtheorem{remark}{Remark}
\newtheorem{observation}{Observation}
\newtheorem{experiment}{Experiment}

%% Custom Commands (Preserved from your original)
\newcommand{\reals}{\mathbb{R}}
\newcommand{\intg}{\mathbb{Z}}
\newcommand{\ind}{\mathbb{I}}
\newcommand{\natr}{\mathbb{N}}
\newcommand{\mc}{\mathcal}
\newcommand{\mb}{\mathbf}
\newcommand{\mbb}{\mathbb}
\newcommand{\mscr}{\mathscr}
\renewcommand{\v}[1]{\boldsymbol{#1}}
\newcommand{\pp}{\mathbb{P}}
\newcommand{\dd}{\mathrm{d}}
\newcommand{\ee}{\mathbb{E}}
\newcommand{\cA}{\mathcal{A}}
\newcommand{\cC}{\mathcal{C}}
\newcommand{\cB}{B}
\newcommand{\cV}{\mathcal{V}}
\newcommand{\poi}{\textsf{Poi}}
\newcommand{\hppp}{\mathscr{P}}
\newcommand{\argmin}{\operatornamewithlimits{argmin}}
\newcommand{\pnon}{\mathcal{P}_n(\lambda)}
\newcommand{\pnonl}{\mathcal{P}_n(n)}
\newcommand{\mx}{\v{x}}
\newcommand{\mX}{\v{X}}
\newcommand{\mL}{\mb{L}}
\newcommand{\mG}{\mb{G}}
\newcommand{\scG}{\mathscr{G}}
\newcommand{\lt}{\left}
\newcommand{\rt}{\right}
\newcommand{\wt}{\widetilde}
\newcommand{\wh}{\widehat}
\newcommand{\wb}{\overline}
\newcommand{\var}{\mathbb{V}\mathrm{ar}}
\newcommand{\vol}{\mathsf{Vol}}
\newcommand{\maxd}{\mathsf{MD}}
\newcommand{\mcc}{\mathsf{MCC}}
\newcommand{\ntg}{\mathsf{NTG}}
\newcommand{\ec}{\mathsf{EC}}
\newcommand{\mcs}{\mathsf{MCS}}
\newcommand{\sref}[2]{\hyperref[#2]{#1~\ref*{#2}}}

\newcommand{\R}{\mathbb{R}}
\newcommand{\Q}{\mathbb{Q}}
\newcommand{\C}{\mathbb{C}}
\newcommand{\N}{\mathbb{N}}
\newcommand{\F}{\mathbb{F}}
\newcommand{\PP}{\mathbb{P}}
\newcommand{\T}{\mathbb{T}}
\newcommand{\Z}{\mathbb{Z}}
\newcommand{\B}{\mathfrak{B}}
\newcommand{\BB}{\mathcal{B}}
\newcommand{\M}{\mathfrak{M}}
\newcommand{\X}{\mathfrak{X}}
\newcommand{\Y}{\mathfrak{Y}}
\newcommand{\E}{\mathbb{E}}
\newcommand{\cP}{\mathcal{P}}
\newcommand{\cS}{\mathcal{S}}
\newcommand{\cZ}{\mathcal{Z}}
\newcommand{\Ltwo}{\mc{L}^2}
\newcommand{\msf}{\mathsf}
\newcommand{\aset}{\mathscr{A}}
\newcommand{\unif}{\mathsf{Unif}}
\newcommand{\mS}{\mb{S}}
\newcommand{\mY}{\mb{Y}}
\newcommand{\mZ}{\mb{Z}}
\newcommand{\mU}{\mb{U}}
\newcommand{\ma}{\mb{a}}
\newcommand{\ms}{\mb{s}}
\newcommand{\my}{\mb{y}}
\newcommand{\mz}{\mb{z}}
\newcommand{\mXi}{\boldsymbol{\Xi}}
\newcommand{\bad}{\mathsf{Bad}}
\newcommand{\ebad}{\mathsf{EBad}}
\newcommand{\res}{\mathsf{Res}}
\newcommand{\bern}{\mathsf{Bern}}
\newcommand{\badvtx}{\mathsf{BadVtx}}
\newcommand{\vbl}{\mathrm{vbl}}

%% Algorithm Formatting (Preserved)
\algnewcommand{\Inputs}[1]{%
  \State \textbf{Inputs:}
  \Statex \hspace*{\algorithmicindent}\parbox[t]{.8\linewidth}{\raggedright #1}
}
\algnewcommand{\Initialize}[1]{%
  \State \textbf{Initialize:}
  \Statex \hspace*{\algorithmicindent}\parbox[t]{.8\linewidth}{\raggedright #1}
}

%% Title and Authors
\title{Uniform Sampling of Proper Graph Colorings via Soft Coloring and Partial Rejection Sampling}
\author[1]{Sarat Moka\thanks{Corresponding author: \href{mailto:s.moka@unsw.edu.au}{s.moka@unsw.edu.au}}}
\author[2]{Ava Vahedi}
\affil[1]{School of Mathematics and Statistics, University of New South Wales, Sydney, Australia}
\affil[2]{Institute of Algebra, Dresden University of Technology, Dresden, Germany}

\begin{document}

\maketitle

\begin{abstract}
We present a new algorithm for the exact uniform sampling of proper \(k\)-colorings of a graph on \(n\) vertices with maximum degree~\(\Delta\). The algorithm is based on partial rejection sampling (PRS) and introduces a soft relaxation of the proper coloring constraint that is progressively tightened until an exact sample is obtained. Unlike coupling from the past (CFTP), the method is inherently parallelizable. We propose a hybrid variant that decomposes the global sampling problem into independent subproblems of size \(O(\log n)\), each solved by any existing exact sampler. This decomposition acts as a {\em complexity reducer}: it replaces the input size~\(n\) with \(O(\log n)\) in the component solver's runtime, so that any improvement in direct methods automatically yields a stronger result. Using an existing CFTP method as the component solver, this improves upon the best known exact sampling runtime for \(k>3\Delta\). Recursive application of the hybrid drives the runtime to \(O(L^{\log^* n}\cdot n\Delta)\), where \(L\) is the number of relaxation levels. We conjecture that \(L\) is bounded independently of~\(n\), which would yield a linear-time parallelizable algorithm for general graphs. Our simulations strongly support this conjecture.
\end{abstract}
\textbf{Keywords:} Exact Sampling, Coupling From The Past, Lov\'{a}sz Local Lemma, Site Percolation, Parallel Algorithms

\section{Introduction}
\label{chp:intro}

Given an undirected graph \(G=(V,E) \) with vertices $V$ and edges $E$, a {\em proper} \(k\)-coloring is an assignment \(\mx:V\rightarrow [k]\), mapping each vertex to a `color' in $[k] = \{1, 2, \dots, k \}$ such that we have \(\mx(v)\neq\mx(w) \) for all vertices \(v, w\in V \) that have an edge \((v,w)\in E \). That is, for the coloring to be proper, no adjacent vertices may share the same color. A sampling algorithm is called {\em perfect} if it generates a sample from a given distribution within finite time. In this paper, the distribution of interest is the uniform distribution on the set of all proper \(k\)-colorings of the graph.

The leading method for generating perfect samples is {\em Coupling From The Past (CFTP)}, first proposed by \cite{PW96}. Some further advancements of CFTP include those by \cite{MH96}, \cite{JSS21}, and \cite{BC20}. However, a disadvantage of CFTP is that it is sequential, and thus cannot take advantage of parallel computing when implemented.

A different framework for exact sampling, called {\em partial rejection sampling} (PRS), was proposed by \cite{GJL17}, inspired by the resampling algorithm of \cite{MT10}. PRS begins by sampling all variables independently from a reference distribution, then iteratively identifies and resamples only a subset of `bad' variables until an acceptable configuration is reached. A key advantage over CFTP is that PRS is inherently parallelizable, since independent components of the resampling set can be processed concurrently. However, when PRS is applied directly to graph coloring, the resampling set necessarily covers the entire graph and the method degenerates into na\"{i}ve rejection sampling \cite{GJL17,BC20}.

In this paper, we overcome this obstacle by introducing {\em \(\gamma\)-soft coloring}. By augmenting each vertex with a continuous auxiliary random variable \(U_v\sim\unif(0,1)\), we create what we call {\em passive states} that prevent the resampling set from expanding to the full graph. This yields a practical PRS-based algorithm for uniformly sampling proper graph colorings. In particular, our contributions are as follows:

\begin{enumerate}
    \item {\bf \(\gamma\)-soft coloring:} By augmenting each vertex with an auxiliary uniform random variable \(U_v\sim\unif(0,1)\), we introduce what we call passive states, which prevent the resampling set from covering the entire graph. This enables PRS to be applied to graph coloring for the first time.
    \item {\bf Inherent parallelizability:} The resampling set decomposes into independent connected components that can be processed concurrently. Unlike CFTP, which is inherently sequential, our algorithm offers a natural avenue for parallel implementation.
    \item {\bf Hybrid algorithm:} We propose a hybrid variant that uses PRS for the global decomposition into small components, then solves each component by an existing exact sampler such as CFTP with bounding chains; see, e.g., \cite{MH96, BC20}. This combines the parallelizability of PRS with the efficiency of CFTP on small subgraphs.
    \item {\bf Complexity reduction:} We show that the \(\gamma\)-soft decomposition acts as a {\em complexity reducer}: it replaces the input size \(n\) in the component solver's runtime with \(O(\log n)\). More precisely, if the component solver runs in time \(T(m)\) on a graph of \(m\) vertices, the hybrid runs in time \(O(n\cdot T(O(\log n))/\log n)\) per \(\gamma\)-level (\sref{Corollary}{cor:general-solver}). Any future improvement in exact sampling for graph coloring automatically yields, via the hybrid, a strictly faster algorithm.
    \item {\bf Improved complexity over the state of the art:} Applied to the CFTP method of \cite{BC20} (\(k>3\Delta\)), the hybrid achieves an expected runtime of \(O\!\lt(L\cdot n(\log\log n)^2\cdot\Delta^2\log\Delta\log k\rt)\), improving upon the \(O(n\log^2 n\cdot\Delta^2\log\Delta\log k)\) runtime of \cite{BC20} applied directly by a factor of \(\log^2 n / (L\cdot(\log\log n)^2)\), where \(L\) is the number of \(\gamma\)-levels.
    \item {\bf Recursive nesting and the path to linear time:} Since the hybrid is itself an exact sampler, it can serve as its own component solver. Repeated nesting replaces \(n\) with successively iterated logarithms, yielding a runtime of \(O(L^{\log^* n}\cdot n\Delta)\) at full recursion depth.\footnote{The {\em iterated logarithm} \(\log^* n\) is the number of times the logarithm must be applied to~\(n\) before the result is at most~\(1\).} Linear-time exact sampling has been achieved on restricted graph classes by \cite{FGY22}, and concurrently claimed for general graphs (for \(k > 3.637\Delta+1\)) by \cite{BH25} using a sequential randomness recycler. We pose the open question of whether \(L\) remains bounded independently of~\(n\); an affirmative answer would yield a linear-time {\em parallelizable} exact sampler for general graphs. Our simulations strongly support this conjecture.
    \item {\bf Simulations and software:} We provide extensive simulations on cycle graphs, grid graphs, complete graphs, and random regular graphs, validating the theoretical predictions. A Python package implementing all the proposed algorithms is publicly available at \url{https://github.com/saratmoka/parkol}.
\end{enumerate}

\autoref{tab:comparison} summarizes the landscape of exact sampling algorithms for uniform \(k\)-colorings and places our contribution in context.

\begin{table}[ht]
\centering\footnotesize
\resizebox{\textwidth}{!}{%
\begin{tabular}{l l l c l}
\hline
Method & Sequential time & Colors & Parallel & Parallel time\\
\hline
\cite{MH96} & \(O(n\log n)\) & \(k\geq\Delta(\Delta+2)\) & No & ---\\
\cite{BC20} & \(O(n\log^2 n)\) & \(k>3\Delta\) & No & ---\\
\cite{JSS21} & \(O(n\log^2 n)\) & \(k\geq\tfrac{8}{3}\Delta+O(\sqrt{\Delta\log\Delta})\) & No & ---\\
\rowcolor{lightteal} \cite{FGY22} & \(O(n)\) & \(k\geq C\Delta\)\(^\dagger\) & No & ---\\
\hline
{\bf Hybrid + \cite{BC20}} & \(O(Ln(\log\log n)^2)\) & \(k>3\Delta\) & {\bf Yes} & \(O(Ln(\log\log n)^2/M)\)\\
{\bf Hybrid (recursive)} & \(O(L^{\log^* n}\cdot n)\) & \(k>3\Delta\) & {\bf Yes} & \(O(L^{\log^* n}\cdot n/M)\)\\
\rowcolor{lightteal} {\bf Hybrid + \cite{FGY22}} & \(O(Ln)\) & \(k\geq C\Delta\)\(^\dagger\) & {\bf Yes} & \(O(L\log n)\)\(^\ddagger\)\\
\hline
\end{tabular}}
\caption{Exact samplers for uniform \(k\)-colorings (\(n\) vertices, max degree~\(\Delta\)). ``Hybrid'' refers to our hybrid \(\gamma\)-PRS (\sref{Algorithm}{alg:hybrid}), combined with the indicated component solver. Runtimes suppress factors depending only on \(\Delta\) and~\(k\). \(L\) is the number of soft-coloring levels (empirically \({\leq}\,20\), conjectured \(O(1)\)). \(M\) is the number of available processors. \(^\dagger\)\cite{FGY22} applies only to graphs with sub-exponential neighborhood growth (e.g.\ lattices~\(\intg^d\)) under strong spatial mixing; all other methods apply to general graphs. \(^\ddagger\)With \(O(n/\log n)\) processors; see \sref{Remark}{rem:lattice-parallel}.}
\label{tab:comparison}
\end{table}

The remaining paper is organized as follows: \sref{Section}{sec:prelim} introduces notation and the graph coloring problem. In \sref{Section}{sec:prs}, we present the PRS framework of \cite{GJL17}. In \sref{Section}{sec:prs-gamma}, we propose \(\gamma\)-soft coloring and present the hybrid algorithm with its parallelization strategy. In \sref{Section}{sec:runtime}, we analyze the runtime complexity, derive conditions for non-degeneration, and prove the asymptotic improvement over direct CFTP. In \sref{Section}{sec:linear-time}, we discuss the implications for linear-time exact sampling and pose the central open problem. \sref{Section}{sec:sim-results} presents simulation results.

\section{Preliminaries}\label{sec:prelim}

We first introduce some notation used throughout the paper. We denote by \(\unif[k]\) the discrete uniform distribution on \([k] = \{1,\ldots,k\}\) and by \(\unif(0,1)\) the continuous uniform distribution on \((0,1)\); we write \(X\sim\mu\) if the distribution of a random object \(X\) is~\(\mu\). For two probability measures \(\mu_1\) and \(\mu_2\) on the same measurable space, \(\mu_1\ll\mu_2\) means that \(\mu_1\) is absolutely continuous with respect to~\(\mu_2\). We write \(e=\exp(1)\).

Let \(G=(V,E)\) be an undirected graph with \(n=|V|\) vertices, \(|E|\) edges, and maximum degree~\(\Delta\). For each vertex~\(v\), let \(N(v)=\{w\in V:(v,w)\in E\}\) denote its set of neighbors and \(d_v=|N(v)|\) its degree. The graph is called \(\Delta\)-{\em regular} if \(d_v=\Delta\) for every \(v\in V\). A \(k\)-coloring of \(G\) is an assignment \(\mx:V\to[k]\). The reference distribution \(\rho\) is the product measure on \([k]^V\) under which every \(k\)-coloring is equally likely. That is, if we associate with each vertex \(v\) an independent random variable \(X_v\sim\unif[k]\), and write \(\mX=(X_v)_{v\in V}\) for the joint process, then $\mX$ has distribution $\rho$.

A \(k\)-coloring \(\mx\) is {\em proper} if \(\mx(v)\neq\mx(w)\) for every edge \((v,w)\in E\). The set of proper colorings is \(\mc{A}=\{\mx\in[k]^V:\mx(v)\neq\mx(w)\;\forall\,(v,w)\in E\}\), and the target distribution \(\mu\) is the uniform distribution on \(\mc{A}\), i.e., \(\mu(\mx)=\rho(\mx\mid\mx\in\mc{A})\). Our goal is to draw exact samples from \(\mu\). When PRS is applied directly to this problem, the resampling set necessarily covers the entire graph and PRS degenerates into na\"{i}ve rejection sampling; refer to  \cite{GJL17,BC20}. This motivates the introduction of \(\gamma\)-soft coloring in \sref{Section}{sec:prs-gamma}.

\section{Partial Rejection Sampling}\label{sec:prs}

In this section we present the partial rejection sampling (PRS) framework introduced by \cite{GJL17}; see also \cite{FGY24} for a comprehensive survey. We follow closely the formulation in \cite{MK20}.

Let \(\mX = \{X_1, X_2, \ldots, X_n\}\) be a set of independent random objects, where each \(X_i\) takes values in some space~\(\mc X\). Write \(\rho\) for the product distribution of~\(\mX\); we call \(\rho\) the {\em reference distribution}. Let \(\{B_v : v \in V_D\}\) be a set of \(m=|V_D|\) {\em bad events}, indexed by elements of a finite set~\(V_D\). Each bad event \(B_v\) depends on a subset of the random objects; let \(\mc I(v) \subseteq \{1,\ldots,n\}\) be the index set such that \(B_v\) depends on \(\{X_i : i \in \mc I(v)\}\) and is independent of the remaining objects. For any \(W \subseteq V_D\), define \(\mc I(W) = \bigcup_{v \in W} \mc I(v)\).

Two bad events \(B_u\) and \(B_v\) are called {\em dependent} if \(\mc I(u) \cap \mc I(v) \neq \emptyset\), i.e., they share at least one random object. The {\em dependency graph} has vertex set~\(V_D\) and an edge between \(u\) and~\(v\) whenever \(B_u\) and \(B_v\) are dependent. For a realization \(\mx = (x_1,\ldots,x_n)\) of~\(\mX\) and a subset \(W \subseteq V_D\), the {\em partial realization} of \(\mx\) restricted to~\(W\) is \(\mx|_W := \{x_i : i \in \mc I(W)\}\). A realization \(\mx'\) is called an {\em extension} of the partial realization \(\mx|_W\) if \(\mx'|_W = \mx|_W\), that is, \(\mx'\) agrees with \(\mx\) on the objects indexed by~\(\mc I(W)\) but may differ elsewhere. We say that a bad event \(B_v\) is {\em disjoint} from \(\mx|_W\) if either \(\mc I(v) \cap \mc I(W) = \emptyset\), or \(B_v\) does not occur under any extension of~\(\mx|_W\).

Let \(\bad(\mx) = \{v \in V_D : \mx \in B_v\}\) be the set of bad events that occur under~\(\mx\), and let \(\mc A = \{\mx : \bad(\mx) = \emptyset\}\) be the {\em acceptable set}. The goal of PRS is to draw exact samples from the {\em target distribution} \(\mu = \rho(\cdot \mid \mc A)\).

\sref{Algorithm}{alg:prs} presents the PRS method. In each iteration, it constructs a {\em resampling set} \(\res \subseteq V_D\) by starting from \(\bad(\mx)\) and expanding through the dependency graph: boundary events that are {\em not} disjoint from \(\mx|_{\res}\) are added to \(\res\), while those that {\em are} disjoint are placed in~\(N\) and the expansion halts at their boundary. Once \(\res\) is determined, all random objects with indices in \(\mc I(\res)\) are resampled from~\(\rho\). For PRS to be effective, disjoint events must exist at the boundary; otherwise, \(\res\) covers all of~\(V_D\) and PRS reduces to na\"{i}ve rejection sampling.

\begin{algorithm}[H]\caption{Partial Rejection Sampling \cite[Algorithm~6]{GJL17}; see also \cite{MK20}}\label{alg:prs}
Draw independent samples \(X_1,\ldots,X_n\) from \(\rho\). Set \(\mx \leftarrow (X_1,\ldots,X_n)\).\\
\While{\(\bad(\mx) \neq \emptyset\)}{
    \(\res \leftarrow \bad(\mx)\),\quad \(N \leftarrow \emptyset\)\\
    \While{\(\partial \res \setminus N \neq \emptyset\)}{
        Let \(D = \{v \in \partial \res \setminus N : B_v \text{ is disjoint from } \mx|_{\res}\}\)\\
        \(N \leftarrow N \cup D\),\quad \(\res \leftarrow \res \cup (\partial \res \setminus N)\)
    }
    Resample the objects \(\{X_i : i \in \mc I(\res)\}\)
}
Output \(\mx\).
\end{algorithm}

Here \(\partial \res\) denotes the boundary of~\(\res\) in the dependency graph, i.e., the set of events adjacent to~\(\res\) but not in~\(\res\). See \cite[Theorem~4.5]{GJL17} for a proof that \sref{Algorithm}{alg:prs} outputs samples from~\(\mu\).

\begin{example}[Hard-Core Model]\label{exp:sshc}\normalfont
In the {\em hard-core model}, each vertex of an undirected graph is independently occupied with probability \(\frac{\lambda}{1+\lambda}\) for some fugacity \(\lambda>0\), and the target distribution is conditioned on no edge having both endpoints occupied. This can be viewed as a \(2\)-coloring problem: each vertex is colored red (occupied) or green (unoccupied), and the constraint forbids adjacent vertices from both being red, while adjacent green vertices are permitted. A bad event \(B_v\) is associated with each vertex~\(v\): it occurs when \(v\) is red and has at least one red neighbor, with \(\mc I(v) = \{v\}\cup N(v)\). Crucially, a green (unoccupied) vertex is always disjoint from any partial realization, since it cannot become bad regardless of its neighbors' configuration. This provides the disjoint events needed for PRS to be effective, and the resampling set can remain strictly smaller than~\(V_D\). Under appropriate conditions on~\(\lambda\) and the maximum degree~\(\Delta\), PRS achieves \(O(n)\) expected runtime for this model; see \cite[Theorem~6.5]{GJL17} and \cite{MK20}.

The proper \(k\)-coloring problem is strictly harder in this regard: since every color can conflict with a neighbor's color, {\em no} vertex state is unconditionally disjoint. As a result, when PRS is applied directly to uniform \(k\)-coloring, the resampling set necessarily covers all of~\(V_D\) and PRS reduces to na\"{i}ve rejection sampling \cite{BC20}. Overcoming this obstacle is the main contribution of the present paper.
\hfill\(\lozenge\)
\end{example}
\section{Perfect Sampling for Graph Colorings}\label{sec:prs-gamma}

In this section, we introduce the \(\gamma\)-soft coloring framework that enables PRS for graph coloring, present the main algorithm and its recursive and hybrid variants, and discuss parallelization.

\subsection{\texorpdfstring{\(\gamma\)}{gamma}-Soft Coloring}\label{sec:soft-coloring}

Consider a random process \(\mX=(X_v)_{v\in V}=(C_v,U_v)_{v\in V}\) where the color \(C_v\) of vertex~\(v\) is an independent random variable from the discrete uniform distribution \(\unif[k]\) and \(U_v\) is an independent random variable from the continuous uniform distribution \(\unif(0,1)\). Then realizations \(\mx=(c_v,u_v)_{v\in V}\) of \(\mX\) are elements of \(([k]\times \left(0,1\right))^V \), that is to say, they are functions taking each vertex \(v\in V\) to both a color and a value between \(0\) and \(1\). We are then interested in defining our reference and target measures on \(([k]\times \left(0,1\right))^V \) as well as a series of `intermediate' measures.

The reference and target measures need only be defined with reference to the color of each vertex: the reference measure \(\rho\) is that where each \(k\)-coloring has equal probability of occurring, while for the target measure \(\mu\), each proper coloring has equal probability of occurring and improper colorings have probability \(0\) of occurring. To get from \(\rho\) to \(\mu\), we however make use of those intermediate measures which are defined at particular \(\gamma\), for \(\gamma\in[0,1]\). Let 
\[
n_v(\gamma, \mx) = |\{w \in V | c_w = c_v\And u_w > \gamma^{d_w} \And (v,w)\in E\}|
\] 
be the number of neighbors of a vertex \(v\) which have the same color as \(v\) and have \(u_w>\gamma^{d_w}\), where \(d_w\) denotes the degree of vertex \(w\) (with the convention \(0^0=1\)). We will sometimes write \(n_v\) instead of \(n_v(\gamma,\mx)\) for brevity. Note that since \(u_v \in (0,1)\), a vertex with \(u_v>\gamma^{n_v}\) necessarily has a same-color neighbor: if it has no such neighbor, then \(n_v=0\) and \(\gamma^{n_v}=1>u_v\).
With this understanding, we associate with each vertex \(v\) a bad event
\[
B_{\gamma,v} = \{u_v>\gamma^{n_v}\},
\]
and define the set of bad vertices as
\[
\bad(\mx, \gamma)=\{v\in V\big|B_{\gamma,v}\text{ occurs}\} = \{v\in V\big|u_v>\gamma^{n_v}\}.
\]
Note that bad events are indexed by the vertex set~\(V\) of the graph, so \(V_D = V\) in the notation of \sref{Section}{sec:prs}. Each bad event \(B_{\gamma,v}\) depends on the variables \((c_w,u_w)\) for \(w\in\{v\}\cup N(v)\), giving \(\mc I(v) = \{v\}\cup N(v)\).

These definitions allow us to introduce what we call a {\em passive state}: a vertex \(v\) is passive at level \(\gamma\) if \(u_v\leq\gamma^{d_v}\). Such a vertex cannot be bad regardless of its neighbors' configuration, because even if all \(d_v\) neighbors share the same color as~\(v\) (so that \(n_v=d_v\)), we still have \(u_v\leq\gamma^{d_v}=\gamma^{n_v}\). In the language of \sref{Section}{sec:prs}, a passive vertex is always disjoint from any partial realization \(\mx|_{\res}\): it cannot become bad no matter how the resampling set is resampled. Consequently, the expansion of the resampling set in \sref{Algorithm}{alg:prs} halts at passive vertices. The existence of passive states is the key property that prevents the resampling set from covering the entire graph, and it is precisely what the auxiliary uniform random variables provide.

Now based on our definition of bad vertices, we can define the acceptable set at each \(\gamma\) as 
\[
\mc{A}_\gamma= \{\mx\, \big|\, \bad(\mx, \gamma) = \emptyset\}  = \{\mx\big|u_v\leq\gamma^{n_v}\,\, \forall\, v\in V \},
\] 
so that a realization is acceptable at that \(\gamma\) if there is no vertex with \(u_v>\gamma^{n_v}\), which necessarily implies that there is no edge with both endpoints the same color and both \(u_v>\gamma^{n_v}\) and \(u_w>\gamma^{n_w}\). Then we can define the intermediate measure at some \(\gamma\) as 
\begin{align}
\label{eqn:eta_defn}
    \eta_\gamma(\mx)=\rho(\mx|\mc{A}_\gamma).
\end{align} 
We call a sample \(\mx\in\mc{A}_\gamma\) a {\em sample of \(\gamma\)-soft coloring}. Note that clearly \(\mc{A}_0=\mc A\), so that \(\eta_0=\mu\). We also have \(\eta_1=\rho\). Importantly, for \(\gamma<\gamma'\), we have \(\eta_\gamma\ll\eta_{\gamma'}\) since $\mc{A}_{\gamma} \subseteq \mc{A}_{\gamma'}$.

The principle of PRS for proper graph coloring with \(\gamma\)-soft coloring acts by applying \sref{Algorithm}{alg:prs} with target measure \(\eta_\gamma\) to a sample \(\mx\) to get as output a sample from~\(\eta_\gamma\). If this sample is not a proper graph coloring, the procedure is repeated at a lower value of \(\gamma\), and so on.
To show that this procedure will result in samples with distribution \(\mu\), we have the following result.

\begin{theorem}
    \label{thm:conv-soft-coloring}
    For any graph, the distribution of $\gamma$-soft coloring converges to the uniform distribution on proper colorings as $\gamma$ goes to zero. That is, for any sample~\(\mx\), 
    \[
    \mu(\mx) = \lim_{\gamma \to 0^+} \eta_\gamma(\mx), \quad \text{for all }\, \, \mx.
    \]
\end{theorem}

\begin{proof}
For \(\mx\notin\mc A\), there exists an edge \((v,w)\) with \(c_v=c_w\). For sufficiently small \(\gamma>0\), this edge forces at least one of \(v,w\) into \(\bad(\mx,\gamma)\), so \(\mx\notin\mc A_\gamma\) and hence \(\eta_\gamma(\mx)=0=\mu(\mx)\).

Now fix \(\mx\in\mc A\). From the definition~\eqref{eqn:eta_defn},
\[
\eta_\gamma (\mx) = \frac{\rho(\mx)}{\rho(\mc A_{\gamma})}.
\]
For any \(\gamma < \gamma'\), we have \(\mc A \subseteq \mc A_{\gamma} \subseteq \mc A_{\gamma'}\), so the sets \(\{\mc A_\gamma\}_{\gamma>0}\) are decreasing as \(\gamma\downarrow0\) and satisfy \(\bigcap_{\gamma>0}\mc A_\gamma = \mc A\). Since \(\rho\) is a probability measure and \(\rho(\mc A_1)<\infty\), continuity of measure from above \cite[Chapter 2]{SB95} gives
\[
\lim_{\gamma \to 0^+} \rho(\mc A_\gamma) =  \rho\lt(\bigcap_{\gamma>0} \mc A_\gamma\rt) = \rho(\mc A).
\]
Therefore,
\[
\lim_{\gamma \to 0^+} \eta_\gamma(\mx)  = \frac{\rho(\mx)}{\rho(\mc A)} = \mu(\mx), \quad \text{for all}\,\, \mx.
\]
\end{proof}

\subsection{The New Algorithm}
We now present our novel algorithm for generating perfect samples of target distribution $\mu$, the distribution of uniformly selected proper graph coloring. In practice, we decrease \(\gamma\) as the algorithm progresses. In particular,  we refer to a decreasing sequence $\{\gamma_\ell \in (0, 1) : \ell \in \mathbb{N}_0\}$ as {\em valid $\gamma$-sequence} if $\gamma_0 = 1 > \gamma_1 > \gamma_2 > \cdots $ and $\gamma_\ell \to 0$ as $\ell \to \infty$. One such valid sequence is given by $\gamma_\ell = 0.9^\ell$, which is used in our simulation results presented in \sref{Subsection}{sec:sim-results}.

\sref{Algorithm}{alg:pc-prs} uniformly samples a proper graph coloring via partial rejection sampling of \(\gamma\)-soft coloring, which we will call \(\gamma\)-PRS. Later in \sref{Section}{sec:sampling-gamma-soft}, we provide a recursive implementation of \(\gamma\)-PRS.

\begin{algorithm}[H]
	\caption{Proper Coloring through PRS}\label{alg:pc-prs}
	Draw a sample \(\mx\) from the reference distribution \(\rho \).\\
    Choose a valid $\gamma$-sequence $\{\gamma_\ell : \ell \geq 0\}$.\\
    Set \(\ell=0\).\\
    \While{there exists an edge \((v,w)\in E\) with \(c_v=c_w\)}{
        \While{$\bad(\mx, \gamma_\ell) \neq \emptyset$}{
            Construct the resampling set \(R\) by expanding from \(\bad(\mx,\gamma_\ell)\) through non-passive vertices, including the passive boundary (see steps (i)--(iii) below).\\
            \(\mx(R) \leftarrow \gamma\text{-PRS}(R, \mx(R), \ell)\) \tcp*{\sref{Algorithm}{alg:rec-prs}}
        }
        \(\ell=\ell+1\)
	}
    Output \(\mx \).
\end{algorithm}
\ \\
\sref{Algorithm}{alg:pc-prs} starts with a sample from the reference distribution $\rho$. For a fixed valid $\gamma$-sequence, starting with $\ell = 0$, it increases $\ell$ by one iteratively. 
For each $\ell$, the inner while loop generates perfect $\gamma_\ell$-soft coloring.
For this, we identify the resampling set~\(R\) and call \(\gamma\)-PRS\((R, \mx(R), \ell)\), which generates a sample of \(\gamma_\ell\)-soft coloring taking the current state \(\mx(R)\) on~\(R\) as the initial realization.

Since bad events are indexed by vertices, the resampling set~\(R\) is a set of vertices. Concretely, \(R\) is constructed as follows:
\begin{enumerate}[label=(\roman*)]
\item Initialise \(R=\bad(\mx,\gamma)\), the set of bad vertices.
\item Expand from~\(R\): for each neighbor \(w\) of \(R\) not yet visited, if \(w\) is non-passive (\(u_w>\gamma^{d_w}\)), add \(w\) to~\(R\) and continue expanding; if \(w\) is passive (\(u_w\leq\gamma^{d_w}\)), mark \(w\) as boundary and do not expand further.
\item Add the passive boundary vertices to~\(R\).
\end{enumerate}
In the notation of \sref{Section}{sec:prs}, the inner set (steps (i)--(ii), excluding the passive boundary) corresponds to~\(\res\), while the full set~\(R\) (including the passive boundary) corresponds to~\(\mc I(\res)\), i.e., the variables that \sref{Algorithm}{alg:prs} resamples. Resampling~\(R\) means drawing fresh values \((c_v,u_v)\sim\unif[k]\times\unif(0,1)\) independently for every \(v\in R\), while keeping all variables outside~\(R\) unchanged.

\begin{theorem}
\label{thm:prs_coloring}
\sref{Algorithm}{alg:pc-prs} halts in finite time almost surely and its output is a uniformly selected proper coloring.
\end{theorem}
\begin{proof}
We first verify that the \(\gamma\)-soft coloring setup satisfies the assumptions of PRS \cite[Theorem 4.5]{GJL17}: the reference distribution \(\rho\) is a product measure on \(([k]\times(0,1))^V\), and the bad events \(\{B_v:v\in V\}\) are determined by the random variables \((C_v,U_v)_{v\in V}\). These conditions are satisfied by construction.

At each level \(\ell\), the inner loop of \sref{Algorithm}{alg:pc-prs} applies PRS with reference distribution \(\rho\) and bad events defined at \(\gamma_\ell\). By \cite[Theorem 4.5]{GJL17}, the output is an exact sample from \(\eta_{\gamma_\ell}\). Denote the configuration after the \(\ell\)-th level by \(\mX^{(\ell)}\), so that \(\mX^{(\ell)} \sim \eta_{\gamma_\ell}\).

Let $T$ denote the first level at which the outer loop terminates, i.e., $T=\inf\{\ell:\mX^{(\ell)}\in\mc A\}$. We show that $\pp(T < \infty) = 1$. Since the events $\{T \leq \ell\}$ are monotonically increasing in $\ell$, by continuity of probability \cite[Chapter 2]{SB95},
\[
\pp(T < \infty) = \lim_{\ell \to \infty} \pp\lt( T \leq \ell \rt).
\]
Further,
\[
\pp\lt( T \leq \ell \rt) \geq \pp\lt( \mX^{(\ell)} \in \mc A\rt) = \eta_{\gamma_\ell} (\mc A).
\]
Since $\gamma_\ell \to 0$ and $\eta_{\gamma_\ell}(\mc A) = \rho(\mc A)/\rho(\mc A_{\gamma_\ell}) \to 1$ by \sref{Theorem}{thm:conv-soft-coloring}, we obtain $\pp(T < \infty) = 1$.

It remains to show that $\mX^{(T)} \sim \mu$. At the terminating level $T$, we have $\mX^{(T)} \sim \eta_{\gamma_T}$ and $\mX^{(T)} \in \mc A$ (by definition of $T$). Therefore the output has distribution $\eta_{\gamma_T}(\cdot \mid \mc A)$. Since $\mc A \subseteq \mc A_{\gamma_T}$ and $\eta_{\gamma_T} = \rho(\cdot \mid \mc A_{\gamma_T})$, for any $\mx \in \mc A$ we have
\[
\eta_{\gamma_T}(\mx \mid \mc A) = \frac{\eta_{\gamma_T}(\mx)}{\eta_{\gamma_T}(\mc A)} = \frac{\rho(\mx)/\rho(\mc A_{\gamma_T})}{\rho(\mc A)/\rho(\mc A_{\gamma_T})} = \frac{\rho(\mx)}{\rho(\mc A)} = \mu(\mx).
\]
Hence $\mX^{(T)} \sim \mu$.
\end{proof}

\subsection{Sampling of \texorpdfstring{\(\gamma \)}{TEXT}-Soft Coloring}
\label{sec:sampling-gamma-soft}

For \sref{Algorithm}{alg:pc-prs} to execute correctly, we require an implementation of $\gamma$-PRS($G, \mx, \ell$) on any graph $G$ and for any $\gamma_\ell$, starting with a realization $\mx$ from $\gamma_{\ell-1}$-soft coloring (i.e., from $\eta_{\gamma_{\ell-1}}$). Since $\eta_0 = \rho$ and each $\eta_\gamma$ is absolutely continuous with respect to $\rho$, one could use any existing algorithm to generate samples from $\eta_{\gamma_\ell}$, taking $\rho$ as the reference measure. Here we provide a recursive algorithm for generating samples of $\gamma_\ell$-soft coloring using samples from $\gamma_j$-soft colorings for $j = 0, 1, \dots, \ell -1$ (\sref{Algorithm}{alg:rec-prs}). In \sref{Subsection}{sec:hybrid}, we demonstrate how existing exact sampling methods such as CFTP can be used in place of this recursion.

\begin{algorithm}[H]
	\caption{$\gamma$-PRS($G, \mx, \ell$)}\label{alg:rec-prs}
    \While{$\bad(\mx, \gamma_\ell) \neq \emptyset$}{
        Construct \(R\) by expanding from \(\bad(\mx,\gamma_\ell)\) through non-passive vertices, including the passive boundary.\\
        Update $\mx$ by resampling all the vertices of $R$ under $\rho$.\\
        Let \(G_1,..., G_a\) be connected components of \(R\).\\
        \For{$i = 1, \dots, a$}{
            \For{\(j = 0, 1,...,\ell\)}{
                \(\mx(G_i) \leftarrow \gamma\text{-PRS}(G_i, \mx(G_i), j)\) \tcp*{recurse}
            }
        }
    }
    Output \(\mx \).
\end{algorithm}

\sref{Algorithm}{alg:rec-prs} is a recursive algorithm whose output are samples from the intermediate measures \(\eta_{\gamma_\ell}\). At each iteration, the resampling set~\(R\) is constructed as in \sref{Algorithm}{alg:pc-prs}, and all vertices in~\(R\) are resampled. Here the connected components \(G_1,\ldots,G_a\) of~\(R\) are the components of the subgraph of~\(G\) induced by the vertices in~\(R\). This is done recursively through the levels until a sample from \(\eta_\ell\) is reached.

The reason why the resampling set is split into connected components is parallelization. This allows the problem to be split into multiple sub-problems, which can then be executed concurrently on different processors, resulting in a reduction in running time.

\subsection{Parallelization}\label{sec:parallel}

A distinctive advantage of PRS over CFTP-based methods is its natural parallelizability. At each iteration of the inner while loop, the resampling set \(R\) decomposes into connected components \(G_1,\ldots,G_a\). Since the components are conditionally independent (each depends only on its own vertices and the passive boundary), they can be processed concurrently on~\(a\) processors with no inter-process communication.

The parallel cost of each iteration is determined by the {\em largest} component: if the components have sizes \(s_1,\ldots,s_a\), the sequential cost is \(\sum_i s_i = |R|\), while the parallel cost (with sufficiently many processors) is \(\max_i s_i\). The speedup is thus \(|R|/\max_i s_i\), which is significant when~\(R\) consists of many small components.

For a fixed \(\gamma\)-sequence, whether the resampling set decomposes into multiple components depends on the graph size and the current \(\gamma\)-value. \autoref{tab:components} reports the average component structure (over 20 independent random colorings) at selected \(\gamma\)-values, for random \(3\)-regular graphs, random \(4\)-regular graphs, and grid graphs.

Several phenomena are evident. First, the number of components is maximized in an intermediate window of \(\gamma\)-values: for \(\gamma\) too close to~\(1\) there are few bad vertices and hence few components, while for \(\gamma\) too small the components merge into one giant component (the percolation transition of \sref{Subsection}{sec:non-degen}). For \(3\)-regular graphs with \(k=15\), this window is approximately \(\gamma\in[0.87,0.93]\); for \(4\)-regular graphs with \(k=20\), it is narrower, around \(\gamma\in[0.93,0.96]\).

Second, the number of components {\em grows with~\(n\)}: at \(\gamma=0.91\), the average number of components increases from~\(3\) at \(n=1000\) to~\(6\) at \(n=2000\) and~\(16\) at \(n=5000\). This is consistent with the percolation theory, which predicts that the sub-critical regime (many small components) becomes more pronounced as \(n\to\infty\).

\begin{table}[ht]
\centering\small
\begin{tabular}{l r r | r r r r r}
\hline
Graph & \(n\) & \(\gamma\) & avg \(|\bad|\) & avg \(|R|\) & avg \(\#\)comp & max \(\#\)comp & avg max comp \\
\hline
\multicolumn{8}{l}{\em Random \(3\)-regular, \(k=15\)} \\
\(n\!=\!1000\) & 1000 & 0.93 & 2.8 & 17 & 2.0 & 6 & 9.1 \\
\(n\!=\!1000\) & 1000 & 0.91 & 4.6 & 37 & 3.4 & 6 & 17.0 \\
\(n\!=\!1000\) & 1000 & 0.89 & 6.6 & 56 & 4.0 & 7 & 27.4 \\
\(n\!=\!2000\) & 2000 & 0.93 & 5.7 & 37 & 4.2 & 7 & 12.1 \\
\(n\!=\!2000\) & 2000 & 0.91 & 8.3 & 65 & 6.0 & 11 & 21.9 \\
\(n\!=\!2000\) & 2000 & 0.89 & 13.8 & 118 & 7.7 & 12 & 40.6 \\
\(n\!=\!5000\) & 5000 & 0.93 & 12.7 & 82 & 9.8 & 18 & 15.4 \\
\(n\!=\!5000\) & 5000 & 0.91 & 22.9 & 170 & 15.5 & 19 & 37.0 \\
\(n\!=\!5000\) & 5000 & 0.89 & 30.9 & 272 & 17.6 & 24 & 55.9 \\
\hline
\multicolumn{8}{l}{\em Grid, \(k=20\)} \\
\(30\!\times\!30\) & 900 & 0.93 & 2.8 & 37 & 2.0 & 5 & 22.6 \\
\(30\!\times\!30\) & 900 & 0.91 & 4.8 & 72 & 3.4 & 7 & 35.1 \\
\(30\!\times\!30\) & 900 & 0.89 & 6.0 & 111 & 3.8 & 7 & 51.8 \\
\(50\!\times\!50\) & 2500 & 0.93 & 8.1 & 91 & 6.5 & 13 & 24.3 \\
\(50\!\times\!50\) & 2500 & 0.91 & 13.2 & 202 & 9.8 & 15 & 43.0 \\
\(50\!\times\!50\) & 2500 & 0.89 & 20.8 & 381 & 11.6 & 18 & 94.8 \\
\hline
\multicolumn{8}{l}{\em Random \(4\)-regular, \(k=20\)} \\
\(n\!=\!1000\) & 1000 & 0.96 & 1.4 & 17 & 0.8 & 2 & 13.6 \\
\(n\!=\!1000\) & 1000 & 0.95 & 1.6 & 18 & 1.1 & 3 & 13.7 \\
\(n\!=\!1000\) & 1000 & 0.93 & 3.5 & 60 & 1.4 & 3 & 52.5 \\
\(n\!=\!2000\) & 2000 & 0.96 & 2.4 & 30 & 1.9 & 4 & 19.4 \\
\(n\!=\!2000\) & 2000 & 0.95 & 3.9 & 50 & 2.2 & 4 & 33.9 \\
\(n\!=\!2000\) & 2000 & 0.93 & 7.8 & 146 & 2.0 & 5 & 130.1 \\
\hline
\end{tabular}
\caption{Average connected component structure of the resampling set at selected \(\gamma\)-values (20 trials). As \(n\) grows, the number of components increases, confirming that parallelization benefits improve with graph size.}
\label{tab:components}
\end{table}

Third, \(4\)-regular graphs exhibit fewer components and a narrower useful \(\gamma\)-window than \(3\)-regular graphs. This is expected: higher degree means denser connectivity among non-passive vertices, so components merge more easily.

These observations suggest an {\em adaptive} parallelization strategy: rather than following a fixed \(\gamma\)-sequence, decrease \(\gamma\) until the resampling set splits into at most \(M\) components (where \(M\) is the number of available processors), solve the components in parallel, and repeat. As \(n\) grows, the window of \(\gamma\)-values yielding multiple components widens, consistent with the percolation theory of \sref{Subsection}{sec:non-degen}. On graphs with sub-exponential neighborhood growth, the hybrid with CFTP achieves \(O(\log n)\) parallel time (see \sref{Remark}{rem:lattice-parallel}).

\subsection{Hybrid \(\gamma\)-PRS}\label{sec:hybrid}

The recursive \sref{Algorithm}{alg:rec-prs} is theoretically clean, but in practice the recursion tree through levels \(0,1,\ldots,\ell\) on each connected component can become expensive. We now describe a {\em hybrid} variant that replaces the recursive inner loop with any existing exact sampler, thereby decoupling the global decomposition power of PRS from the local sampling problem on each component.

The key observation is the following. At level \(\ell\), after the resampling set \(R\) is constructed and we decompose it into connected components \(G_1,\ldots,G_a\), the components are conditionally independent given the configuration on the passive boundary. It therefore suffices to produce, on each component \(G_i\), a sample from \(\eta_{\gamma_\ell}\) restricted to the vertices of~\(G_i\), with the vertices outside \(G_i\) held fixed. Any exact sampler that targets this conditional distribution may be used.

Since each \(\eta_{\gamma_\ell}\) is simply the reference distribution \(\rho\) conditioned on \(\mc{A}_{\gamma_\ell}\), one natural choice is na\"{i}ve rejection sampling (NRS) on the component: repeatedly resample all vertices of \(G_i\) from \(\rho\) until the \(\gamma_\ell\)-soft constraint is satisfied. Because the components are typically much smaller than the full graph, the acceptance probability of NRS on a component is substantially higher than on the full graph, making this approach practical.

Alternatively, one can use Coupling From The Past (CFTP) with bounding chains on each component. CFTP produces a uniformly selected proper \(k\)-coloring of the induced subgraph~\(G_i\), independently of the configuration outside~\(G_i\). A proper coloring automatically satisfies the \(\gamma\)-soft constraint at every level (since a proper coloring has \(n_v=0\) for all~\(v\), giving \(\gamma^{n_v}=1>u_v\) always), so the result is always in~\(\mc A_{\gamma_\ell}\). Two CFTP methods are applicable:
\begin{itemize}
\item The bounding chain method of \cite{MH96}, which requires \(k\geq\Delta(\Delta+2)\) for polynomial runtime.
\item The improved method of \cite{BC20}, which requires only \(k>3\Delta\) and runs in expected time \(O\!\lt(n\log^2 n\cdot\Delta^2\log\Delta\log k\rt)\).
\end{itemize}
Because the components are typically small, CFTP coalesces quickly even on subgraphs where the global runtime bound would be pessimistic.

This hybrid approach is formalized in \sref{Algorithm}{alg:hybrid}.

\begin{algorithm}[H]
	\caption{Hybrid \(\gamma\)-PRS}\label{alg:hybrid}
	Draw a sample \(\mx\) from the reference distribution \(\rho\).\\
    Choose a valid \(\gamma\)-sequence \(\{\gamma_\ell:\ell\geq0\}\). Set \(\ell=0\).\\
    \While{there exists an edge \((v,w)\in E\) with \(c_v=c_w\)}{
        \While{\(\bad(\mx,\gamma_\ell)\neq\emptyset\)}{
            Construct \(R\) by expanding from \(\bad(\mx,\gamma_\ell)\) through non-passive vertices, including the passive boundary.\\
            Let \(G_1,\ldots,G_a\) be the connected components of \(R\).\\
            \For{\(i=1,\ldots,a\) {\em (independently, in parallel)}}{
                Sample a uniform proper \(k\)-coloring of the subgraph~\(G_i\).\\
                Draw fresh \(u_v\sim\unif(0,1)\) for each \(v\in G_i\).\\
            }
        }
        \(\ell=\ell+1\).
	}
    Output \(\mx\).
\end{algorithm}

\begin{theorem}\label{thm:hybrid_correct}
Suppose the exact sampler used on each component \(G_i\) produces a uniform proper \(k\)-coloring of the subgraph~\(G_i\), independently of the configuration outside~\(G_i\). Then \sref{Algorithm}{alg:hybrid} halts in finite time almost surely and its output is a uniformly selected proper coloring.
\end{theorem}
\begin{proof}
Let \(\eta'_{\gamma}\) denote the distribution of the configuration \(\mx\) at the end of the inner while loop at level~\(\gamma\) (that is, when \(\bad(\mx,\gamma)=\emptyset\) is first achieved).

{\em Support.}
By construction, \(\eta'_\gamma\) is supported on \(\mc A_\gamma\), since the inner while loop exits only when no vertex is bad at~\(\gamma\). Moreover, every proper coloring \(\mx\in\mc A\) is in~\(\mc A_\gamma\) (since \(n_v=0\) implies \(\gamma^{n_v}=1>u_v\) for all~\(v\)), so \(\mc A\subseteq\mc A_\gamma\).

{\em Uniform on proper colorings.}
We show that \(\eta'_\gamma(\mx)\) is the same for all \(\mx\in\mc A\). Consider the last iteration of the inner while loop that produces the accepted configuration. Let \(R\) be the resampling set and \(G_1,\ldots,G_a\) its connected components. The colors outside~\(R\) are fixed at some configuration~\(\mx_{V\setminus R}\). Each component~\(G_i\) receives an independent uniform proper \(k\)-coloring from the exact sampler, so every proper coloring of~\(G_i\) is equally likely. The configuration is accepted when \(\bad(\mx,\gamma)=\emptyset\). Among all colorings of~\(R\) that are proper on each component, those satisfying the acceptance condition form a set~\(S(\mx_{V\setminus R})\), and each element of~\(S(\mx_{V\setminus R})\) has equal probability (by the uniformity of the component sampler and the independence across components). In particular, for any two proper colorings \(\mx,\mx'\in\mc A\) that agree outside~\(R\), we have \(\eta'_\gamma(\mx)=\eta'_\gamma(\mx')\). Since the argument applies to every possible~\(R\) and external configuration, and since the u-values are drawn independently from \(\unif(0,1)\), the distribution \(\eta'_\gamma\) assigns equal probability to all \(\mx\in\mc A\). That is, \(\eta'_\gamma(\cdot\mid\mc A) = \mu\).

{\em Almost sure halting.}
Let \(T=\inf\{\ell:\mX^{(\ell)}\in\mc A\}\). Since \(\eta'_{\gamma_\ell}\) is supported on~\(\mc A_{\gamma_\ell}\) and \(\mc A\subseteq\mc A_{\gamma_\ell}\), we have \(\eta'_{\gamma_\ell}(\mc A)>0\). As \(\gamma_\ell\to0\), the sets \(\mc A_{\gamma_\ell}\) decrease to~\(\mc A\), so \(\eta'_{\gamma_\ell}(\mc A)\to1\). By continuity of probability, \(\pp(T<\infty)=1\).

{\em Output distribution.}
At level~\(T\), the output satisfies \(\mX^{(T)}\in\mc A\) and is drawn from \(\eta'_{\gamma_T}(\cdot\mid\mc A)=\mu\) by the uniformity argument above.
\end{proof}

\begin{remark}\normalfont
The hybrid approach has two practical advantages over the fully recursive \sref{Algorithm}{alg:rec-prs}:
\begin{enumerate}
\item {\em Reduced complexity.} The recursive algorithm requires \(O(\ell)\) nested calls per level, leading to a recursion tree whose depth grows with the number of levels. The hybrid avoids this by solving each component in a single call to an existing sampler.
\item {\em Parallelizability.} The components \(G_1,\ldots,G_a\) are independent and can be processed concurrently. This is already noted in \sref{Algorithm}{alg:rec-prs}, but the hybrid makes it especially attractive because each component is solved by a self-contained subroutine with no inter-component communication.
\end{enumerate}
Our simulation results (\sref{Subsection}{sec:sim-results}) show that the hybrid with NRS as the component solver reduces the total number of resamplings by up to two orders of magnitude compared to plain iterative PRS. Using CFTP as the component solver yields further improvements, as the BC20 method \cite{BC20} requires only \(k>3\Delta\) and CFTP coalesces quickly on the small components (see \autoref{tab:hybrid-cftp}).
\hfill\(\lozenge\)
\end{remark}

\section{Runtime Analysis}\label{sec:runtime}

We analyze the runtime complexity of \sref{Algorithm}{alg:pc-prs} and the hybrid variant \sref{Algorithm}{alg:hybrid}. We first derive the probability that a vertex is bad at a given~\(\gamma\), which is the key quantity governing convergence of both algorithms.

For clarity, the analysis is carried out for \(\Delta\)-regular graphs. The results extend to arbitrary graphs with maximum degree~\(\Delta\), since the \(\Delta\)-regular case is the worst case: a vertex of degree \(d_v < \Delta\) has a smaller probability of being bad (fewer same-color neighbors) and a higher probability of being passive (\(\gamma^{d_v} > \gamma^{\Delta}\)). Consequently, the resampling set is smaller and the algorithms converge at least as fast as on the corresponding \(\Delta\)-regular graph.

\subsection{Probability of a Bad Vertex}

All probabilities and expectations in this section are with respect to the reference distribution~\(\rho\). Since \(u_v\sim\unif(0,1)\) is independent of \(n_v(\gamma,\mx)\) (which depends on \(c_v\) and on the neighbors' colors and \(u\)-values, but not on~\(u_v\)), we have
\[\mbb P_\rho(u_v>\gamma^{n_v})=\mbb E_\rho\left[\mbb P_\rho(u_v>\gamma^{n_v}\mid n_v)\right]=\mbb E_\rho\left[1-\gamma^{n_v}\right].\]

To evaluate this, note that \(n_v\) counts the neighbors \(w\) of \(v\) with \(c_w=c_v\) and \(u_w>\gamma^{d_w}\). For a neighbor \(w\), let \(F_w\) be the event that \(c_w=c_v\) and \(u_w>\gamma^{d_w}\). Since the colors and \(u\)-values are independent across vertices,
\[\mbb P_\rho(F_w)=\frac{1}{k}(1-\gamma^{d_w}).\]

Although the events \(\{F_w:w\in N(v)\}\) all depend on~\(c_v\), they are mutually independent: since \(c_v,c_w,c_{w'}\) are independent, \(\pp_\rho(c_w=c_v,\,c_{w'}=c_v)=1/k^2=\pp_\rho(c_w=c_v)\,\pp_\rho(c_{w'}=c_v)\), and the \(u\)-values are independent across vertices. Hence \(n_v=\sum_{w\in N(v)}\ind(F_w)\) is a sum of independent Bernoulli random variables. Therefore,
\begin{align}\label{eqn:E_gamma_nv}
\mbb E_\rho\left[\gamma^{n_v}\right]&=\prod_{w\in N(v)}\mbb E_\rho\left[\gamma^{\ind(F_w)}\right]=\prod_{w\in N(v)}\left(1-\frac{(1-\gamma)(1-\gamma^{d_w})}{k}\right),
\end{align}
where we used \(\mbb E_\rho[\gamma^{\ind(F_w)}]=1\cdot\mbb P_\rho(F_w^c)+\gamma\cdot\mbb P_\rho(F_w)=1-(1-\gamma)\mbb P_\rho(F_w)\). Combining,
\begin{align}\label{eqn:P_bad}
\mbb P_\rho(v\in\bad(\mx,\gamma))=\mbb P_\rho(u_v>\gamma^{n_v})=1-\prod_{w\in N(v)}\left(1-\frac{(1-\gamma)(1-\gamma^{d_w})}{k}\right).
\end{align}

For a \(\Delta\)-regular graph, this simplifies to
\begin{align}\label{eqn:P_bad_regular}
\mbb P_\rho(v\in\bad(\mx,\gamma))=1-\left(1-\frac{(1-\gamma)(1-\gamma^{\Delta})}{k}\right)^{\Delta}.
\end{align}
For a general graph with maximum degree~\(\Delta\), this is an upper bound on \(\mbb P_\rho(v\in\bad(\mx,\gamma))\).

Note that as \(\gamma\to1\), \(\mbb P_\rho(v\in\bad)\to0\), confirming that at the reference level \(\gamma_0=1\) there are no bad vertices. As \(\gamma\to0\), \(\mbb P_\rho(v\in\bad)\to1-(1-1/k)^\Delta\), which is the probability that vertex \(v\) shares a color with at least one neighbor under the reference distribution.

Similarly, the probability of a vertex being passive is
\begin{align}\label{eqn:P_passive}
\mbb P_\rho(v\text{ is passive at }\gamma) = \mbb P_\rho(u_v\leq\gamma^{d_v}) = \gamma^{d_v}.
\end{align}
At \(\gamma=1\), every vertex is passive (\(\gamma^{d_v}=1\)), so the resampling set is empty and PRS does nothing. At \(\gamma=0\), no vertex with degree greater than one is passive (\(\gamma^{d_v}=0\) for \(d_v\geq1\)), so the resampling set covers the entire graph and PRS degenerates into na\"{i}ve rejection sampling.

\subsection{Expected Number of Bad Vertices}

By linearity of expectation and \eqref{eqn:P_bad}, the expected number of bad vertices under the reference distribution \(\rho\) is
\begin{align}\label{eqn:E_bad_size}
\mbb E_\rho\left[|\bad(\mx,\gamma)|\right] = \sum_{v\in V}\mbb P_\rho(v\in\bad(\mx,\gamma)) = \sum_{v\in V}\left(1-\prod_{w\in N(v)}\left(1-\frac{(1-\gamma)(1-\gamma^{d_w})}{k}\right)\right).
\end{align}
For a \(\Delta\)-regular graph, this equals
\begin{align}\label{eqn:E_bad_regular}
\mbb E_\rho\left[|\bad(\mx,\gamma)|\right] = n\left(1-\left(1-\frac{(1-\gamma)(1-\gamma^{\Delta})}{k}\right)^{\Delta}\right).
\end{align}
For a general graph with maximum degree~\(\Delta\), the same expression provides an upper bound: \(\mbb E_\rho[|\bad(\mx,\gamma)|] \leq n\!\left(1-\!\left(1-\frac{(1-\gamma)(1-\gamma^{\Delta})}{k}\right)^{\Delta}\right)\!\), since vertices of degree \(d_v < \Delta\) have a smaller probability of being bad.

\subsection{Non-Degeneration Condition: the Percolation Threshold}\label{sec:non-degen}

The resampling set \(R\) is found by expanding from the bad vertices through non-passive vertices (see \sref{Algorithm}{alg:pc-prs}). If the non-passive vertices form a giant connected component in the graph, then even a single bad vertex can cause \(R\) to cover the entire graph. We now derive a condition under which this is avoided.

A vertex \(v\) is non-passive at level~\(\gamma\) with probability \(1-\gamma^{d_v}\). On a \(\Delta\)-regular graph, the non-passive vertices form a random subset where each vertex is included independently with probability
\begin{align}\label{eqn:q_nonpassive}
q(\gamma,\Delta) = 1-\gamma^{\Delta}.
\end{align}
In {\em site percolation}, each vertex of a graph is independently retained with some probability~\(q\) and removed otherwise. The connected cluster of any retained vertex in a graph of maximum degree~\(\Delta\) is stochastically dominated by a Galton--Watson branching process with offspring mean \(q(\Delta-1)\): from any vertex, at most \(\Delta-1\) new neighbors can be explored, each independently retained with probability~\(q\). When \(q(\Delta-1)<1\), the branching process is subcritical and all clusters are finite, with sizes decaying exponentially (see \cite[Section~10.1, Theorem~6.75]{Grimm99}). On a finite graph with \(n\) vertices, this implies that all clusters have size \(O(\log n)\) with high probability. Thus, non-passive vertices do not percolate when
\begin{align}\label{eqn:non-perc}
1 - \gamma^{\Delta} < \frac{1}{\Delta - 1}, \quad\text{equivalently}\quad \gamma > \gamma^* := \left(\frac{\Delta - 2}{\Delta - 1}\right)^{1/\Delta}.
\end{align}
We call \(\gamma^*\) the {\em critical gamma}. For \(\gamma > \gamma^*\), the expansion from any bad vertex reaches only a bounded neighborhood, keeping \(|R|\) proportional to \(|\bad|\) rather than~\(n\).

\begin{proposition}\label{prop:gamma_crit}
For a \(\Delta\)-regular graph with \(\Delta\geq3\), the critical gamma satisfies
\begin{align}
\gamma^* = \left(\frac{\Delta-2}{\Delta-1}\right)^{1/\Delta} \in (0,1),
\end{align}
with \(\gamma^*\to1\) as \(\Delta\to\infty\). In particular, \(\gamma^*\approx0.794\) for \(\Delta=3\), \(\gamma^*\approx0.904\) for \(\Delta=4\), and \(\gamma^*\approx0.944\) for \(\Delta=5\).
\end{proposition}
\begin{proof}
The formula follows from solving \(1-\gamma^\Delta = 1/(\Delta-1)\) for \(\gamma\). Since \(\Delta\geq3\), the ratio \((\Delta-2)/(\Delta-1)\) lies in \((0,1)\), so \(\gamma^*\in(0,1)\). To see that \(\gamma^*\to1\), write \(\gamma^* = \exp\!\lt(\frac{1}{\Delta}\log\frac{\Delta-2}{\Delta-1}\rt)\). Since \(\log\frac{\Delta-2}{\Delta-1}<0\) and \(\frac{1}{\Delta}\log\frac{\Delta-2}{\Delta-1}\to0\) as \(\Delta\to\infty\) (because \(\log\frac{\Delta-2}{\Delta-1}= \log(1-\frac{1}{\Delta-1})\sim -\frac{1}{\Delta-1}\)), we have \(\gamma^*\to1\).
\end{proof}

For the algorithm to avoid degeneration, we need it to find a proper coloring before \(\gamma\) drops below \(\gamma^*\). Since the \(\gamma\)-sequence \(\gamma_\ell=\gamma_{\rm base}^\ell\) decreases geometrically, the number of `effective' levels (above \(\gamma^*\)) is
\begin{align}\label{eqn:ell_star}
\ell^* = \left\lfloor \frac{\log \gamma^*}{\log \gamma_{\rm base}} \right\rfloor.
\end{align}
For \(\gamma_{\rm base}=0.9\) and \(\Delta=4\), this gives \(\ell^*=\lfloor\log(0.904)/\log(0.9)\rfloor=0\), meaning even the first non-trivial level (\(\gamma_1=0.9\)) is near the percolation boundary. This motivates choosing a slower decay rate, such as \(\gamma_\ell = (1-1/(2\Delta))^\ell\), which yields more effective levels.

\subsection{Sufficient Condition on \(k\)}\label{sec:sufficient-k}

For \sref{Algorithm}{alg:pc-prs} and \sref{Algorithm}{alg:hybrid} to terminate efficiently, we need: (i) the resampling set remains small at each level above \(\gamma^*\), and (ii) PRS converges at each such level. We now derive a sufficient condition on \(k\) that ensures this.

\begin{theorem}[Non-degeneration condition]\label{thm:non-degen}
Let \(G\) be a \(\Delta\)-regular graph on \(n\) vertices with \(\Delta\geq3\), and let \(k\) be the number of colors. Define
\begin{align}
\alpha(\gamma) := \frac{(1-\gamma)(1-\gamma^{\Delta})}{k}.
\end{align}
If
\begin{align}\label{eqn:k_sufficient}
k \geq e\,\Delta^3\,(1-\gamma^*)(1-(\gamma^*)^{\Delta}),
\end{align}
where \(\gamma^*\) is as in \sref{Proposition}{prop:gamma_crit}, then for all \(\gamma\geq\gamma^*\):
\begin{enumerate}[label=(\roman*)]
    \item the expected number of bad vertices satisfies \(\mbb E_\rho[|\bad(\mx,\gamma)|]\leq n\,\Delta\,\alpha(\gamma)\);
    \item the non-passive vertices do not percolate, and the expected size of the resampling set is \(\mbb E_\rho[|\res(\mx,\gamma)|] = O\!\lt(\mbb E_\rho[|\bad(\mx,\gamma)|]\rt)\);
    \item the Lov\'asz Local Lemma (LLL) condition \(e\cdot\mbb P_\rho(v\in\bad)\cdot (\Delta^2+1) \leq 1\) is satisfied, ensuring that \(\mbb P_\rho(\bad(\mx,\gamma)=\emptyset)>0\) and hence that PRS terminates almost surely.
\end{enumerate}
\end{theorem}

\begin{proof}
(i) By Bernoulli's inequality, \(1-(1-y)^{\Delta}\leq \Delta y\) for \(y\in[0,1]\) and \(\Delta\geq1\). Applying this with \(y=\alpha(\gamma)\in[0,1]\) and \eqref{eqn:P_bad_regular},
\begin{align*}
\mbb P_\rho(v\in\bad(\mx,\gamma)) = 1-\left(1-\alpha(\gamma)\right)^{\Delta} \leq \Delta\,\alpha(\gamma).
\end{align*}
By linearity of expectation, \(\mbb E_\rho[|\bad|]\leq n\,\Delta\,\alpha(\gamma)\).

(ii) By the definition of \(\gamma^*\), for \(\gamma\geq\gamma^*\) the non-passive fraction \(q = 1-\gamma^{\Delta}\leq 1/(\Delta-1)\), so \(q(\Delta-1)\leq 1\). As argued in \sref{Subsection}{sec:non-degen}, the cluster of any non-passive vertex is dominated by a subcritical Galton--Watson process with offspring mean \(q(\Delta-1)<1\). The expected cluster size is therefore bounded by \(1/(1-q(\Delta-1))\), a constant depending on \(\gamma\) and~\(\Delta\) but not on~\(n\). Since each bad vertex lies in such a cluster, the total resampling set satisfies \(\mbb E_\rho[|\res|] = O(\mbb E_\rho[|\bad|])\).

(iii) The dependency graph of the bad events \(\{B_v:v\in V\}\) has maximum degree at most \(\Delta^2\): two events \(B_v\) and \(B_w\) are dependent whenever \(\mc I(v)\cap\mc I(w)\neq\emptyset\), which requires that \(v\) and \(w\) are neighbors or share a common neighbor, since \(\mc I(v)=\{v\}\cup N(v)\). The number of such vertices~\(w\) is at most \(\Delta+\Delta(\Delta-1)=\Delta^2\), since \(v\) has \(\Delta\) neighbors and each neighbor has at most \(\Delta-1\) other neighbors. By the symmetric form of the Lov\'asz Local Lemma \cite[Theorem 5.1.1]{AS16}, if \(e\cdot\mbb P_\rho(v\in\bad)\cdot(\Delta^2+1)\leq1\) then \(\mbb P_\rho(\bad=\emptyset)>0\). Using the bound from~(i), this requires \(e\,\Delta\,\alpha(\gamma)\cdot(\Delta^2+1)\leq1\), which is implied by \(\alpha(\gamma)\leq 1/(e\,\Delta^3)\). At \(\gamma=\gamma^*\), this becomes \((1-\gamma^*)(1-(\gamma^*)^{\Delta})/k\leq 1/(e\,\Delta^3)\), which is precisely condition~\eqref{eqn:k_sufficient}. Since \(\mbb P_\rho(\bad=\emptyset)>0\) and each PRS iteration resamples the variables in~\(R\) independently from~\(\rho\), the algorithm terminates almost surely; correctness follows from \cite[Theorem~4.5]{GJL17}.
\end{proof}

\begin{remark}\normalfont
The condition~\eqref{eqn:k_sufficient} applies to \sref{Algorithm}{alg:pc-prs} (plain PRS), which requires the LLL contraction in~(iii). For the hybrid \sref{Algorithm}{alg:hybrid}, only parts~(i) and~(ii) are needed; the hybrid's runtime is analyzed separately in \sref{Subsection}{sec:hybrid-runtime} under a weaker condition.

Evaluating \eqref{eqn:k_sufficient} numerically (see \autoref{tab:k-bounds}), the bound is remarkably mild. Since \((1-\gamma^*)(1-(\gamma^*)^{\Delta})\to0\) as \(\Delta\to\infty\) (because \(\gamma^*\to1\)), the right-hand side of~\eqref{eqn:k_sufficient} converges to \(e/(\Delta-1)\approx 2.718/(\Delta-1)\) for large~\(\Delta\). In particular, for \(\Delta\geq5\), the bound is less than~\(\Delta\) and is therefore automatically satisfied whenever \(k>\Delta\) (which is needed for a proper coloring to exist). The bound imposes an additional constraint only for small~\(\Delta\): for \(\Delta=3\) we need \(k\geq8\), and for \(\Delta=4\) we need \(k\geq6\).

\begin{table}[ht]
\centering
\begin{tabular}{c|cccccccc}
\(\Delta\) & 3 & 4 & 5 & 6 & 7 & 8 & 10 & 15 \\
\hline
\(\gamma^*\) & 0.794 & 0.904 & 0.944 & 0.964 & 0.974 & 0.981 & 0.988 & 0.995\\
\(k\) bound in \eqref{eqn:k_sufficient} & 7.6 & 5.6 & 4.7 & 4.3 & 4.0 & 3.8 & 3.5 & 3.2\\
\end{tabular}
\caption{Lower bound on \(k\) from \sref{Theorem}{thm:non-degen} for various \(\Delta\).}
\label{tab:k-bounds}
\end{table}
\hfill\(\lozenge\)
\end{remark}

The bound in \sref{Theorem}{thm:non-degen} guarantees non-degeneration at levels above the percolation threshold. Below \(\gamma^*\), the resampling set may cover the entire graph. However, at this stage the \(\gamma\)-soft constraint has already eliminated most improper edges: our simulations (\sref{Subsection}{sec:sim-results}) confirm that for \(k/\Delta\geq3\) the algorithm typically finds a proper coloring within the effective levels, although the number of PRS iterations per level increases as \(k/\Delta\) decreases toward this threshold.

\subsection{Runtime of the Hybrid Algorithm}\label{sec:hybrid-runtime}

The hybrid \(\gamma\)-PRS (\sref{Algorithm}{alg:hybrid}) replaces the recursive inner loop with an exact sampler on each connected component of the resampling set. We first analyze the per-level cost with NRS as the component solver, then show that replacing NRS with the CFTP method of \cite{BC20} yields an asymptotically faster algorithm than applying \cite{BC20} directly to the full graph.

\begin{lemma}[NRS acceptance on a component]\label{lem:nrs-accept}
Let \(H\) be a connected subgraph of a \(\Delta\)-regular graph with \(|H|=s\) vertices. After resampling all vertices of \(H\) from~\(\rho\), the probability that the \(\gamma\)-soft constraint is satisfied on~\(H\) (i.e., \(\bad(\mx,\gamma)\cap H=\emptyset\)) is at least
\[
\mbb P_\rho(\text{accept}) \geq 1 - s\,\Delta\,\alpha(\gamma),
\]
where \(\alpha(\gamma) = (1-\gamma)(1-\gamma^\Delta)/k\) as before.
\end{lemma}
\begin{proof}
By \eqref{eqn:P_bad_regular} and Bernoulli's inequality (\(1-(1-y)^\Delta\leq\Delta y\) for \(y\in[0,1]\)), each vertex \(v\in H\) has \(\mbb P_\rho(v\in\bad)\leq\Delta\,\alpha(\gamma)\). The result follows from a union bound: \(\mbb P_\rho(\bad\cap H\neq\emptyset)\leq s\,\Delta\,\alpha(\gamma)\).
\end{proof}

The expected number of NRS trials for a component of size \(s\) is therefore at most \(1/(1-s\,\Delta\,\alpha)\), provided \(s\,\Delta\,\alpha < 1\).

\begin{remark}[Inner loop iterations with CFTP]\normalfont\label{rem:cftp-iterations}
When CFTP is used as the component solver (\sref{Algorithm}{alg:hybrid}), each call produces a proper \(k\)-coloring of the component independently of the external configuration (\sref{Theorem}{thm:hybrid_correct}). Since the coloring within each component is proper, no bad vertex can arise from edges {\em within} a component. However, a vertex on the boundary of a component may share its new color with a neighbor {\em outside} the component, potentially creating a bad vertex at level~\(\gamma_\ell\). In that case, the inner while loop of \sref{Algorithm}{alg:hybrid} iterates: a new resampling set is constructed around the newly bad vertices, and fresh CFTP calls resolve them.

The probability of a cross-edge conflict at a single boundary vertex is at most~\(1/k\) (the probability that CFTP assigns the same color as the external neighbor). Since the passive boundary has size \(O(\mbb E_\rho[|\bad|])\) when \(q < p_c\), the expected number of new bad vertices per iteration is \(O(\mbb E_\rho[|\bad|]/k)\), which is small for large~\(k\). In practice, we observe that the inner while loop terminates within very few iterations (typically~\(1\)--\(3\)).
\hfill\(\lozenge\)
\end{remark}

To bound the total per-level cost, we use the cluster size distribution when \(q(\Delta-1)<1\). The connected cluster of any non-passive vertex is stochastically dominated by a Galton--Watson branching process: from any vertex, at most \(\Delta-1\) neighbors can extend the cluster, each independently non-passive with probability~\(q\). The offspring distribution is therefore dominated by \(\mathrm{Bin}(\Delta-1,q)\) with mean \(q(\Delta-1)<1\), so the process is subcritical. The probability that the cluster has size \(\geq s\) decays exponentially:
\begin{align}\label{eqn:cluster-tail}
\mbb P_\rho(|C_v|\geq s) \leq \exp(-c\,s), \quad\text{where}\quad c=c(\gamma,\Delta)=\log\frac{1}{q(\Delta-1)} > 0,
\end{align}
and the expected cluster size is \(1/(1-q(\Delta-1))\).

\begin{theorem}[Per-level runtime of the hybrid]\label{thm:hybrid-runtime}
Let \(G\) be a \(\Delta\)-regular graph on \(n\) vertices with \(\Delta\geq3\). Consider the hybrid \(\gamma\)-PRS (\sref{Algorithm}{alg:hybrid}) with NRS as the component solver at a level with \(\gamma > \gamma^*\). If
\begin{align}\label{eqn:k_hybrid}
\Delta\,\alpha(\gamma) < c(\gamma,\Delta), \quad\text{equivalently}\quad k > \frac{\Delta\,(1-\gamma)(1-\gamma^\Delta)}{c(\gamma,\Delta)},
\end{align}
where \(c(\gamma,\Delta)\) is the percolation decay rate~\eqref{eqn:cluster-tail}, then the expected total NRS cost at this level is \(O(n)\).
\end{theorem}
\begin{proof}
The resampling set \(R\) decomposes into connected components \(G_1,\ldots,G_a\), each contained in a subcritical percolation cluster. By \sref{Lemma}{lem:nrs-accept}, the NRS acceptance probability on a component of size~\(s_i\) is at least \(1-s_i\,\Delta\,\alpha\), so the expected number of NRS trials is at most \(1/(1-s_i\,\Delta\,\alpha)\). Since each trial costs \(O(s_i)\), the total NRS cost across all components is at most
\[
\sum_{i=1}^a \frac{s_i}{1-s_i\,\Delta\,\alpha}.
\]
When \(q(\Delta-1) < 1\), each component has size \(s_i = O(\log n)\) with high probability (\sref{Lemma}{lem:max-comp}). For a component contained in a percolation cluster of size~\(s\), the contribution to the total cost is \(s/(1-s\,\Delta\,\alpha)\). Taking expectations and using the cluster size tail bound~\eqref{eqn:cluster-tail}, the expected contribution of a single bad vertex is
\[
\mbb E_\rho\lt[\frac{|C|}{1-|C|\,\Delta\,\alpha}\rt].
\]
Since \(\mbb P_\rho(|C|\geq s)\leq\exp(-cs)\), this expectation is finite whenever \(\Delta\,\alpha < c\), which is ensured by~\eqref{eqn:k_hybrid}. Under this condition, the expected total NRS cost is bounded by \(\mbb E_\rho[|\bad|]\cdot M\), where \(M<\infty\) is a constant depending on \(\gamma,\Delta,k\) but not on~\(n\). Since \(\mbb E_\rho[|\bad|]\leq n\,\Delta\,\alpha\) by \sref{Theorem}{thm:non-degen}(i), the expected total NRS cost per level is \(O(n)\).
\end{proof}

\begin{remark}\normalfont
The condition~\eqref{eqn:k_hybrid} is strictly weaker than the LLL condition~\eqref{eqn:k_sufficient} of \sref{Theorem}{thm:non-degen}. For example, at \(\gamma=0.95\) on a \(4\)-regular graph, the percolation decay rate is \(c\approx0.59\), and~\eqref{eqn:k_hybrid} requires only \(k\geq 1\) (trivially satisfied), whereas the LLL condition at \(\gamma^*\approx0.90\) requires \(k\geq 6\). More generally, at levels well above \(\gamma^*\) the hybrid requires \(k\) only slightly larger than~\(\Delta\), compared to \(k=O(\Delta^3)\) for the worst-case LLL condition at \(\gamma\to0\). This improvement arises because the hybrid exploits the small size of subcritical clusters rather than requiring a global contraction of the bad set.
\hfill\(\lozenge\)
\end{remark}

\subsection{Asymptotic Improvement and Comparison with Existing Methods}\label{sec:hybrid-vs-direct}

When CFTP is used as the component solver in the hybrid, we can compare the total runtime against applying CFTP directly to the full graph. The key observation is that in subcritical percolation, the largest component has size \(O(\log n)\), so the CFTP cost on each component replaces \(\log^2 n\) with \((\log\log n)^2\).

\begin{lemma}[Maximum component size]\label{lem:max-comp}
Let \(G\) be a \(\Delta\)-regular graph on \(n\) vertices with \(\Delta\geq3\), and let \(\gamma>\gamma^*\). Then the connected components of the resampling set \(R\) each have size at most \(O(\log n)\) with high probability.
\end{lemma}
\begin{proof}
The resampling set \(R\) is contained in the union of subcritical percolation clusters seeded at bad vertices. In subcritical site percolation on a \(\Delta\)-regular graph at occupation probability \(q=1-\gamma^\Delta < 1/(\Delta-1)\), the probability that any vertex belongs to a cluster of size \(\geq s\) is at most \(\exp(-cs)\) where \(c = \log(1/(q(\Delta-1)))>0\) (by the branching process bound~\eqref{eqn:cluster-tail}). By a union bound over all \(n\) vertices, the probability that any cluster exceeds size \(s = (1+\varepsilon)\log n / c\) is at most \(n\cdot\exp(-cs) = n^{-\varepsilon}\to 0\).
\end{proof}

\begin{theorem}[Hybrid with BC20 vs.\ direct BC20]\label{thm:hybrid-vs-direct}
Let \(G\) be a \(\Delta\)-regular graph on \(n\) vertices with \(\Delta\geq3\) and \(k>3\Delta\). Let \(L\) denote the number of \(\gamma\)-levels used by \sref{Algorithm}{alg:hybrid}. Then:
\begin{enumerate}[label=(\roman*)]
\item The expected cost of applying BC20 \cite{BC20} directly to \(G\) is
\[
T_{\mathrm{direct}} = O\!\lt(n\log^2 n\cdot\Delta^2\log\Delta\log k\rt).
\]
\item At each level \(\ell\) with \(\gamma_\ell>\gamma^*\), the expected cost of running BC20 on all components of the resampling set is
\[
T_{\mathrm{level}} = O\!\lt(n\,(\log\log n)^2\cdot\Delta^2\log\Delta\log k\rt).
\]
\item The expected total cost of the hybrid over \(L\) levels is
\[
T_{\mathrm{hybrid}} = O\!\lt(L\cdot n\,(\log\log n)^2\cdot\Delta^2\log\Delta\log k + L\cdot n\Delta\rt),
\]
which is asymptotically faster than \(T_{\mathrm{direct}}\) whenever \(L = o\!\lt(\log^2 n / (\log\log n)^2\rt)\).
\end{enumerate}
\end{theorem}

\begin{proof}
Part~(i) is \cite[Theorem 1.1]{BC20}.

For part~(ii), at a level with \(\gamma>\gamma^*\), the resampling set \(R\) decomposes into components \(G_1,\ldots,G_a\) with sizes \(s_1,\ldots,s_a\). By \sref{Lemma}{lem:max-comp}, \(s_{\max}:=\max_i s_i = O(\log n)\) with high probability. The BC20 cost on a component of size \(s_i\) is \(O(s_i\log^2 s_i\cdot\Delta^2\log\Delta\log k)\). Summing over all components,
\begin{align*}
\sum_{i=1}^a s_i\,\log^2 s_i &\leq \lt(\sum_{i=1}^a s_i\rt)\cdot \log^2 s_{\max} \leq |R|\cdot\log^2(O(\log n)) = O\!\lt(n\,(\log\log n)^2\rt),
\end{align*}
where we used \(|R|\leq n\) and \(\log^2(C\log n) = O((\log\log n)^2)\). Adding the \(O(n\Delta)\) cost of computing \(\bad\) and \(R\) gives part~(ii).

Part~(iii) follows by summing over \(L\) levels. Comparing with~(i), \(T_{\mathrm{hybrid}} < T_{\mathrm{direct}}\) whenever \(L\cdot(\log\log n)^2 < \log^2 n\), i.e., \(L < \log^2 n/(\log\log n)^2\). In our simulations, \(L\) ranges from~3 to~20, so the condition is amply satisfied for any practical \(n\).
\end{proof}

\sref{Theorem}{thm:hybrid-vs-direct} is in fact a special case of a more general principle. The \(\gamma\)-soft decomposition reduces the problem size from \(n\) to \(O(\log n)\); any improvement in the component solver automatically propagates to the hybrid.

\begin{corollary}[General component solver]\label{cor:general-solver}
Suppose there exists an exact sampling algorithm \(\mc S\) for uniform proper \(k\)-colorings that, on a graph with \(m\) vertices and maximum degree~\(\Delta\), runs in expected time \(T_{\mc S}(m,\Delta,k)\). Then the hybrid \(\gamma\)-PRS with \(\mc S\) as the component solver has expected total runtime
\[
O\!\lt(L\cdot n\cdot\frac{T_{\mc S}(O(\log n),\,\Delta,\,k)}{O(\log n)} + L\cdot n\Delta\rt),
\]
where \(L\) is the number of \(\gamma\)-levels and the \(T_{\mc S}(O(\log n),\Delta,k)/O(\log n)\) term is the amortised per-vertex cost of running \(\mc S\) on components of size \(O(\log n)\).
\end{corollary}
\begin{proof}
By \sref{Lemma}{lem:max-comp}, each component has size \(s_i = O(\log n)\) with high probability. The cost of running \(\mc S\) on all components at one level is \(\sum_i T_{\mc S}(s_i,\Delta,k) \leq \sum_i s_i \cdot \frac{T_{\mc S}(s_{\max},\Delta,k)}{s_{\max}}\), since \(T_{\mc S}(m,\Delta,k)/m\) is non-decreasing in \(m\) for any reasonable algorithm. Since \(\sum_i s_i \leq n\) and \(s_{\max} = O(\log n)\), the per-level cost is \(n\cdot T_{\mc S}(O(\log n),\Delta,k)/O(\log n)\). Summing over \(L\) levels and adding the \(O(n\Delta)\) PRS overhead per level gives the result.
\end{proof}

\begin{remark}\normalfont
\sref{Corollary}{cor:general-solver} shows that the \(\gamma\)-soft decomposition acts as a {\em complexity reducer}: it replaces the argument \(n\) in the component solver's runtime with \(O(\log n)\). For any algorithm whose runtime is super-linear in \(n\), this yields a strict improvement. For example:
\begin{itemize}
\item BC20 \citep{BC20} has \(T_{\mc S}(m) = O(m\log^2 m\cdot\Delta^2\log\Delta\log k)\). The hybrid gives \(O(L\cdot n(\log\log n)^2\cdot\Delta^2\log\Delta\log k)\), saving a factor of \(\log^2 n/(L\cdot(\log\log n)^2)\); this is \sref{Theorem}{thm:hybrid-vs-direct}.
\item Huber \citep{MH96} has \(T_{\mc S}(m) = O(m\log m\cdot\frac{k-\Delta}{k-\Delta(\Delta+2)})\). The hybrid gives \(O(L\cdot n\log\log n\cdot\frac{k-\Delta}{k-\Delta(\Delta+2)})\).
\item If a future algorithm achieves \(T_{\mc S}(m) = O(m\,\mathrm{polylog}(m)\cdot f(\Delta,k))\), the hybrid would yield \(O(L\cdot n\,\mathrm{polylog}(\log n)\cdot f(\Delta,k))\).
\end{itemize}
Thus, any future improvement in exact sampling algorithms for graph coloring automatically translates, via the hybrid, into an even faster algorithm through the \(n\to O(\log n)\) reduction.
\hfill\(\lozenge\)
\end{remark}

\begin{remark}[Parallel speedup on lattice graphs]\normalfont\label{rem:lattice-parallel}
On graphs with sub-exponential neighborhood growth (such as lattices~\(\intg^d\)), \cite{FGY22} achieves \(O(n)\) sequential runtime for perfect sampling of colorings when strong spatial mixing holds. Using \cite{FGY22} as the component solver in our hybrid does not improve the sequential cost (\(T_{\mc S}(m)=O(m)\) gives a per-level cost of~\(O(n)\)). However, the decomposition into independent components of size \(O(\log n)\) enables a {\em parallel} speedup: with \(O(n/\log n)\) processors, each component is solved in \(O(\log n)\) time, yielding a total parallel runtime of \(O(L\cdot\log n)\). For bounded~\(L\), this is \(O(\log n)\), which is exponentially faster than the \(O(n)\) sequential runtime of~\cite{FGY22}, which, being based on a sequential Gibbs sampler, cannot be parallelized in this way.
\hfill\(\lozenge\)
\end{remark}

Moreover, since the hybrid itself is an exact sampler, it can serve as its own component solver, leading to a recursive application.

\begin{theorem}[Recursive hybrid]\label{thm:recursive-hybrid}
Let \(T^{(d)}(n)\) denote the expected per-level runtime of the \(d\)-times nested hybrid \(\gamma\)-PRS on an \(n\)-vertex \(\Delta\)-regular graph with \(k>3\Delta\), where
\begin{itemize}
\item \(T^{(0)}(n)\) is the cost of a direct component solver (e.g., BC20), and
\item \(T^{(d)}(n)\) uses the \((d-1)\)-times nested hybrid as the component solver.
\end{itemize}
Write \(\log^{(d)} n\) for the \(d\)-th iterated logarithm, and let \(L\) be the number of \(\gamma\)-levels at each recursion depth. Then
\[
T^{(d)}(n) = O\!\lt(L^d\cdot n\cdot\frac{T^{(0)}\!\lt(O(\log^{(d)} n)\rt)}{O(\log^{(d)} n)} + L^d\cdot n\Delta\rt).
\]
In particular, if \(T^{(0)}(m)=O(m\log^2 m\cdot f(\Delta,k))\) for some function~\(f\), then
\[
T^{(d)}(n) = O\!\lt(L^d\cdot n\cdot\lt(\log^{(d+1)} n\rt)^2\cdot f(\Delta,k) + L^d\cdot n\Delta\rt).
\]
At depth \(d=\log^* n - O(1)\), the iterated logarithm \(\log^{(d)} n = O(1)\), the component solver runs in \(O(f(\Delta,k))\), and the total cost becomes \(O(L^{\log^* n}\cdot n\Delta)\).
\end{theorem}

\begin{proof}
At each recursion depth, \sref{Corollary}{cor:general-solver} replaces the problem size \(m\) with \(O(\log m)\). Starting from \(n\), after \(d\) applications the component size is \(\log^{(d)} n\). The multiplicative overhead \(L^d\) arises because each of the \(d\) recursion depths contributes a factor of~\(L\) levels. The substitution \(T^{(0)}(O(\log^{(d)} n))\) and the cost formula follow from iterated application of \sref{Corollary}{cor:general-solver}. At \(d = \log^* n - O(1)\), we have \(\log^{(d)} n = O(1)\), so the component solver cost is \(O(1)\cdot f(\Delta,k)\), and the dominant term is the PRS decomposition overhead \(O(L^d\cdot n\Delta)\).
\end{proof}

\begin{remark}\normalfont
For practical values of~\(n\), the iterated logarithm \(\log^* n\leq 5\), so the recursion depth is at most~\(5\). If \(L\) is bounded (as observed in our simulations, where \(L\leq 20\)), the overhead \(L^{\log^* n}\) is a moderate constant. In this regime, the recursive hybrid achieves an essentially {\em linear-time} exact sampler: \(O(n\cdot\mathrm{poly}(\Delta,k,L))\). In practice, a single level of nesting (\(d=1\)) already captures the main improvement, and deeper nesting offers diminishing returns.
\hfill\(\lozenge\)
\end{remark}

\subsubsection*{Comparison with existing methods}

Na\"{i}ve rejection sampling (NRS) repeatedly draws a coloring from~\(\rho\) and accepts if it is proper. The acceptance probability satisfies \(\rho(\mc A)\leq((k-1)/k)^{|E|}\), so the expected number of iterations is at least \((k/(k-1))^{|E|}\), which grows exponentially in the number of edges. For a \(\Delta\)-regular graph, this is at least \((k/(k-1))^{n\Delta/2}\).

The leading alternative is coupling from the past (CFTP), introduced by \cite{PW96}. Subsequent improvements by \cite{MH96}, \cite{BC20}, and \cite{JSS21} have progressively reduced both the runtime and the minimum number of colors required. The runtimes and color conditions for these methods are summarized in \autoref{tab:comparison} in the introduction. All CFTP-based methods are inherently sequential, unlike our PRS-based approach which is parallelizable. In our hybrid algorithm (\sref{Algorithm}{alg:hybrid}), we use CFTP methods of \cite{MH96} and \cite{BC20} as component solvers on the small subgraphs produced by the \(\gamma\)-soft decomposition.

\section{Towards Linear-Time Exact Sampling}\label{sec:linear-time}

Most exact samplers for uniform \(k\)-colorings have super-linear runtime in the number of vertices~\(n\). The CFTP method of \cite{MH96} achieves \(O(n\log n\cdot\mathrm{poly}(\Delta,k))\), and \cite{BC20} achieves \(O(n\log^2 n\cdot\mathrm{poly}(\Delta,k))\). Even for approximate sampling via Glauber dynamics, the mixing time is \(\Theta(n\log n)\) due to the coupon-collector barrier. Linear-time exact sampling has been achieved in restricted settings: \cite{GJL17} proved an \(O(n)\) bound for the hard-core model, and \cite{FGY22} achieved \(O(n)\) for colorings on graphs with sub-exponential neighborhood growth (such as lattices~\(\intg^d\)) when strong spatial mixing holds. Very recently, \cite{BH25} claimed \(O(n\Delta)\) runtime on general graphs for \(k > 3.637\Delta + 1\) using a sequential randomness recycler; this is an arXiv preprint and has not yet been peer-reviewed. However, none of these methods are parallelizable.

The recursive hybrid of \sref{Theorem}{thm:recursive-hybrid} brings us within reach of this goal. Recall that with \(d\) levels of nesting, the runtime is
\[
O\!\lt(L^d \cdot n \cdot (\log^{(d+1)} n)^2 \cdot f(\Delta,k)\rt),
\]
where \(f(\Delta,k) = \Delta^2\log\Delta\log k\), \(\log^{(d)}\) denotes the \(d\)-th iterated logarithm, and \(L\) is the number of \(\gamma\)-levels per recursion depth. At depth \(d = \log^* n - O(1)\), the component solver cost vanishes and the runtime reduces to
\begin{align}\label{eqn:recursive-limit}
O(L^{\log^* n}\cdot n\Delta).
\end{align}
Since \(\log^* n \leq 5\) for all practical~\(n\) (indeed, \(\log^* 2^{65536} = 5\)), the factor \(L^{\log^* n}\) is a moderate constant whenever \(L\) is bounded. The entire runtime is then {\em linear in~\(n\)}, up to a multiplicative constant depending on \(\Delta\), \(k\), and~\(L\).

Whether this linear-time guarantee holds depends on a single question:

\begin{quote}
\textbf{Open Problem.} For \(k > C\Delta\) (with \(C\) a suitable constant), does the number of \(\gamma\)-levels \(L\) in \sref{Algorithm}{alg:hybrid} remain bounded as \(n\to\infty\)? That is, does there exist a constant \(L_0 = L_0(\Delta, k)\), independent of~\(n\), such that \(L \leq L_0\) with high probability?
\end{quote}

An affirmative answer would immediately yield, via~\eqref{eqn:recursive-limit}, a {\em linear-time parallelizable} exact sampler for uniform proper \(k\)-colorings on general graphs:
\[
O\!\lt(L_0^{\,\log^*\! n}\cdot n\Delta\rt) = O(n\cdot\mathrm{poly}(\Delta,k)),
\]
since \(L_0^{\,\log^*\! n}\) is a constant for any fixed \(\Delta\) and~\(k\).

We note that proving \(L = O(1)\) requires showing that at the \(\gamma\)-levels above the percolation threshold \(\gamma^*\), the probability of finding a proper coloring is bounded away from zero uniformly in~\(n\). We leave the resolution of this question as an important direction for future work.

An important observation is that the choice of \(\gamma\)-sequence affects the value of~\(L\) but not the underlying difficulty. Since \(\eta_{\gamma_{{\ell}-1}}(\cdot\mid\mc A_{\gamma_\ell})=\eta_{\gamma_\ell}\), a sample from level \(\ell-1\) that already satisfies the constraint at level~\(\ell\) is automatically a sample from~\(\eta_{\gamma_\ell}\), requiring no resampling. We define the {\em effective} number of levels \(L_{\mathrm{eff}}\) as the number of levels at which PRS actually performs resampling (i.e., \(\bad(\mx,\gamma_\ell)\neq\emptyset\)); the remaining levels are ``free'' and require only a check. \autoref{tab:effective-levels} shows that with a fine \(\gamma\)-sequence (\(\gamma_\ell = 0.99^\ell\)), most levels are skipped and \(L_{\mathrm{eff}}\) is consistently small (\(1\)--\(7\)), independent of the step factor. This suggests that the algorithm needs to do real work at only a few critical \(\gamma\)-values, and that \(L_{\mathrm{eff}}\) is bounded.

\begin{table}[ht]
\centering
\begin{tabular}{l r r r | r r | r r | r r}
\hline
 & & & & \multicolumn{2}{c|}{\(\gamma_\ell=0.99^\ell\)} & \multicolumn{2}{c|}{\(\gamma_\ell=0.95^\ell\)} & \multicolumn{2}{c}{\(\gamma_\ell=0.9^\ell\)} \\
Graph & \(n\) & \(\Delta\) & \(k\) & \(L_{\mathrm{eff}}\) & skip & \(L_{\mathrm{eff}}\) & skip & \(L_{\mathrm{eff}}\) & skip \\
\hline
Petersen & 10 & 3 & 5 & 1 & 15 & 1 & 3 & 1 & 2 \\
\(C_{20}\) & 20 & 2 & 5 & 7 & 252 & 11 & 33 & 8 & 15 \\
Grid \(5\!\times\!5\) & 25 & 4 & 10 & 5 & 177 & 5 & 32 & 5 & 14 \\
3-reg \(n\!=\!50\) & 50 & 3 & 10 & 9 & 174 & 9 & 24 & 8 & 8 \\
\(K_{10}\) & 10 & 9 & 15 & 3 & 117 & 3 & 22 & 3 & 10 \\
Grid \(5\!\times\!5\) & 25 & 4 & 20 & 2 & 57 & 2 & 11 & 2 & 5 \\
\hline
\end{tabular}
\caption{Effective levels \(L_{\mathrm{eff}}\) (where PRS resamples) vs.\ skipped levels (where the sample already satisfies the next constraint), for three \(\gamma\)-sequences. \(L_{\mathrm{eff}}\) is consistently small and independent of the step factor.}
\label{tab:effective-levels}
\end{table}

\section{Simulation Results}\label{sec:sim-results}

We implemented all proposed algorithms in Python. The code is available as the open-source package \texttt{parkol},\footnote{\url{https://github.com/saratmoka/parkol}} installable via \texttt{pip install parkol}. All experiments use the valid \(\gamma\)-sequence \(\gamma_\ell=0.9^\ell\). We compare the iterative variant of \(\gamma\)-PRS (\sref{Algorithm}{alg:pc-prs}), the hybrid (\sref{Algorithm}{alg:hybrid}), and na\"{i}ve rejection sampling (NRS) on several graph families.

\autoref{tab:small-graphs} reports results on small graphs where both methods terminate. For these, \(\gamma\)-PRS and NRS are both fast, but \(\gamma\)-PRS already uses fewer total resamplings on denser graphs (e.g., \(K_{10}\) with \(k=15\): 22 resamplings vs.\ 67 NRS iterations).

\begin{table}[ht]
\centering
\begin{tabular}{l r r r | r r | r}
\hline
Graph & \(n\) & \(\Delta\) & \(k\) & Levels & Resamp. & NRS iter. \\
\hline
Cycle \(C_{10}\) & 10 & 2 & 5 & 9 & 5 & 7 \\
Petersen & 10 & 3 & 5 & 3 & 1 & 4 \\
\(K_6\) & 6 & 5 & 10 & 4 & 1 & 4 \\
\(K_{10}\) & 10 & 9 & 15 & 13 & 22 & 67 \\
Cycle \(C_{20}\) & 20 & 2 & 4 & 11 & 24 & 147 \\
\hline
\end{tabular}
\caption{Comparison of \(\gamma\)-PRS (iterative) and NRS on small graphs.}
\label{tab:small-graphs}
\end{table}

\autoref{tab:scaling} demonstrates the effect of the ratio \(k/\Delta\) on performance. When \(k/\Delta\) is large (say \(\geq 5\)), the algorithm converges quickly because a large fraction of vertices are passive at each level, keeping the resampling set small. When \(k/\Delta\) is close to \(1\), the resampling set tends to cover the entire graph and PRS degenerates toward NRS.

\begin{table}[ht]
\centering
\begin{tabular}{l r r r r | r r r}
\hline
Graph & \(n\) & \(\Delta\) & \(k\) & \(k/\Delta\) & Levels & Resamp. & Vtx resamp. \\
\hline
Grid \(5\times5\) & 25 & 4 & 10 & 2.5 & 19 & 218 & 5\,446 \\
Grid \(5\times5\) & 25 & 4 & 20 & 5.0 & 7 & 8 & 200 \\
Grid \(10\times10\) & 100 & 4 & 20 & 5.0 & 11 & 3\,733 & 372\,947 \\
Grid \(10\times10\) & 100 & 4 & 50 & 12.5 & 7 & 10 & 983 \\
3-reg \(n=100\) & 100 & 3 & 20 & 6.7 & 19 & 1\,168 & 116\,490 \\
3-reg \(n=200\) & 200 & 3 & 50 & 16.7 & 12 & 226 & 45\,083 \\
3-reg \(n=200\) & 200 & 3 & 100 & 33.3 & 6 & 2 & 207 \\
\hline
\end{tabular}
\caption{Scaling behavior of \(\gamma\)-PRS (iterative) with varying \(k/\Delta\). Random regular graphs are generated with the indicated degree.}
\label{tab:scaling}
\end{table}

The simulation results confirm two key observations: (i) the number of levels required grows modestly (typically 7--19), while the per-level cost depends heavily on \(k/\Delta\); and (ii) for sufficiently large \(k/\Delta\), the total number of resamplings remains small even as \(n\) grows, consistent with the passive-state fraction \(\gamma^{\Delta}\) remaining above the site percolation threshold for the graph.

\subsubsection*{Hybrid \(\gamma\)-PRS}

\autoref{tab:hybrid} compares plain iterative PRS (\sref{Algorithm}{alg:pc-prs}) with the hybrid variant (\sref{Algorithm}{alg:hybrid}) using NRS as the component solver. The hybrid approach dramatically reduces the total number of resamplings: on the \(5\times5\) grid with \(k=10\), the hybrid uses only~3 resamplings compared to~218 for plain PRS, a reduction of over 70\(\times\). This is because the connected components of the resampling set are small, and NRS on a small component has high acceptance probability.

\begin{table}[ht]
\centering
\begin{tabular}{l r r r | r r | r r}
\hline
 & & & & \multicolumn{2}{c|}{PRS (iterative)} & \multicolumn{2}{c}{Hybrid-NRS} \\
Graph & \(n\) & \(\Delta\) & \(k\) & Levels & Resamp. & Levels & Resamp. \\
\hline
Petersen & 10 & 3 & 5 & 3 & 1 & 3 & 1 \\
\(K_{10}\) & 10 & 9 & 15 & 13 & 22 & 13 & 3 \\
Grid \(5\times5\) & 25 & 4 & 10 & 19 & 218 & 18 & 3 \\
Grid \(10\times10\) & 100 & 4 & 20 & 11 & 3\,733 & 18 & 10 \\
3-reg \(n=50\) & 50 & 3 & 20 & 5 & 1 & 5 & 1 \\
Grid \(10\times10\) & 100 & 4 & 50 & 7 & 10 & 9 & 4 \\
\hline
\end{tabular}
\caption{Comparison of plain iterative PRS with the hybrid variant using NRS as the component solver. The hybrid reduces total resamplings significantly by exploiting the small size of connected components.}
\label{tab:hybrid}
\end{table}

The hybrid variant uses slightly more levels in some cases (because the component solver re-randomizes the entire component rather than making targeted local changes), but the total cost is lower because each resampling step resolves the component in a single call rather than through many PRS iterations.

\subsubsection*{CFTP component solvers}

We compare two CFTP-based component solvers: Huber's bounding chain method \cite{MH96}, which has a polynomial runtime guarantee for \(k\geq\Delta(\Delta+2)\), and the method of Bhandari and Chakraborty \cite{BC20}, which guarantees polynomial runtime for \(k>3\Delta\). Both algorithms are correct (produce exact uniform samples) for any \(k>\Delta\); the bounds above are for the runtime guarantee only.

A natural question is whether Huber's method works in practice below its theoretical threshold. \autoref{tab:hybrid-cftp} compares both CFTP methods at \(k>3\Delta\), including values well below \(\Delta(\Delta+2)\). Remarkably, Huber's method performs well in this regime: on the small components arising from PRS decomposition, it coalesces quickly even without its polynomial guarantee. At \(k>3\Delta\), both methods achieve similar performance, with Huber often slightly faster due to its simpler update rule.

\begin{table}[ht]
\centering
\begin{tabular}{l r r r r | r r | r r}
\hline
 & & & & & \multicolumn{2}{c|}{Hybrid-Huber} & \multicolumn{2}{c}{Hybrid-BC20} \\
Graph & \(n\) & \(\Delta\) & \(k\) & \(\Delta(\Delta\!+\!2)\) & Levels & Resamp. & Levels & Resamp. \\
\hline
\(C_{50}\) & 50 & 2 & 7 & 8 & 10 & 4 & 13 & 4 \\
Grid \(5\times5\) & 25 & 4 & 13 & 24 & 5 & 1 & 5 & 1 \\
3-reg \(n\!=\!50\) & 50 & 3 & 10 & 15 & 6 & 2 & 6 & 2 \\
3-reg \(n\!=\!100\) & 100 & 3 & 10 & 15 & 7 & 3 & 7 & 3 \\
4-reg \(n\!=\!50\) & 50 & 4 & 13 & 24 & 5 & 1 & 5 & 1 \\
\hline
\end{tabular}
\caption{Hybrid \(\gamma\)-PRS at \(3\Delta < k < \Delta(\Delta+2)\): Huber \cite{MH96} vs.\ BC20 \cite{BC20}. Despite operating below its theoretical guarantee \(k\geq\Delta(\Delta+2)\), Huber's method coalesces on the small components produced by PRS and is consistently faster than BC20 due to its simpler update rule.}
\label{tab:hybrid-cftp}
\end{table}

\autoref{tab:hybrid-huber} compares both methods at \(k\geq\Delta(\Delta+2)\).

\begin{table}[ht]
\centering
\begin{tabular}{l r r r | r r | r r}
\hline
 & & & & \multicolumn{2}{c|}{Hybrid-Huber} & \multicolumn{2}{c}{Hybrid-BC20} \\
Graph & \(n\) & \(\Delta\) & \(k\) & Levels & Resamp. & Levels & Resamp. \\
\hline
\(C_{50}\) & 50 & 2 & 8 & 10 & 3 & 17 & 4 \\
3-reg \(n\!=\!50\) & 50 & 3 & 15 & 6 & 2 & 5 & 2 \\
3-reg \(n\!=\!100\) & 100 & 3 & 15 & 9 & 4 & 6 & 4 \\
Grid \(10\times10\) & 100 & 4 & 24 & 4 & 2 & 18 & 3 \\
4-reg \(n\!=\!50\) & 50 & 4 & 24 & 4 & 1 & 4 & 1 \\
\hline
\end{tabular}
\caption{Hybrid \(\gamma\)-PRS at \(k\geq\Delta(\Delta+2)\): Huber \cite{MH96} vs.\ BC20 \cite{BC20}. Both CFTP methods operate within their theoretical guarantees in this regime.}
\label{tab:hybrid-huber}
\end{table}

The key finding is that Huber's simpler bounding chain method works well in practice even below its theoretical threshold \(k\geq\Delta(\Delta+2)\), because the components produced by PRS decomposition are small enough for rapid coalescence. Both methods time out at \(k\leq 2\Delta\) (e.g., Petersen at \(k=5\), Grid at \(k=8\)), and both fail for \(k\leq\Delta\). For the hybrid \(\gamma\)-PRS, we therefore recommend Huber's method as the default component solver when \(k>3\Delta\), with NRS as the fallback for smaller~\(k\).

\section{Conclusion}\label{sec:conclusion}

We introduced \(\gamma\)-soft coloring, a framework that enables partial rejection sampling to be applied to the problem of uniformly sampling proper \(k\)-colorings. The key idea is to augment each vertex with an auxiliary uniform random variable, creating passive states that prevent the resampling set from covering the entire graph. This overcomes a fundamental limitation of existing PRS methods for graph coloring.

Building on this framework, we proposed a hybrid algorithm that decomposes the global sampling problem into independent subproblems on small connected components, each solved by an existing exact sampler such as CFTP. We proved that this decomposition acts as a complexity reducer: it replaces the input size~\(n\) with \(O(\log n)\) in the component solver's runtime (\sref{Corollary}{cor:general-solver}), yielding an asymptotic improvement over all known direct methods (\sref{Theorem}{thm:hybrid-vs-direct}). The hybrid can be applied recursively (\sref{Theorem}{thm:recursive-hybrid}), driving the runtime to \(O(L^{\log^* n}\cdot n\Delta)\).

Two features distinguish our approach from existing CFTP-based methods. First, the algorithm is inherently parallelizable: the independent components can be processed concurrently with no inter-component communication. Second, the framework is modular: any future improvement in exact sampling for graph coloring automatically translates into a faster hybrid algorithm.

An important open question remains: whether the number of \(\gamma\)-levels~\(L\) is bounded independently of~\(n\). An affirmative answer would yield a linear-time parallelizable exact sampler for uniform proper \(k\)-colorings. Our simulations provide strong evidence for this conjecture, with \(L\) remaining between~1 and~20 across all tested graph families and showing no growth with~\(n\).

All algorithms are implemented in the open-source Python package \texttt{parkol}, available at \url{https://github.com/saratmoka/parkol}.

%% Suppress DOI links in references
\makeatletter
\def\doi#1{}
\makeatother

\bibliographystyle{abbrvnat}
\bibliography{Refs}

\appendix
\section{Graph Families Used in Simulations}\label{app:graphs}

We briefly describe each graph family appearing in the simulation results of \sref{Subsection}{sec:sim-results}. In all cases, \(n\) denotes the number of vertices, \(|E|\) the number of edges, and \(\Delta\) the maximum degree.

\begin{description}[style=nextline,leftmargin=2.5cm]

\item[Cycle graph \(C_n\).]
The \(n\) vertices \(\{0,1,\dots,n-1\}\) are arranged in a cycle, with edges \((i,\, i+1\!\!\mod n)\). Every vertex has degree~\(2\), so \(\Delta=2\) and \(|E|=n\). The chromatic number is~\(2\) if \(n\) is even and~\(3\) if \(n\) is odd.

\item[Petersen graph.]
A well-known \(3\)-regular graph on \(n=10\) vertices and \(|E|=15\) edges. It has chromatic number~\(3\), girth~\(5\) (no triangle or \(4\)-cycle), and is vertex-transitive.

\item[Complete graph \(K_n\).]
Every pair of vertices is connected by an edge, giving \(|E|=\binom{n}{2}\) and \(\Delta=n-1\). The chromatic number is~\(n\), so at least \(k=n\) colors are needed for a proper coloring. This is the densest possible simple graph and provides a stress test for the algorithm.

\item[Grid graph \(m\times m\).]
Vertices are placed at the integer lattice points \(\{(i,j):0\leq i,j\leq m-1\}\) with edges between horizontally and vertically adjacent vertices. This gives \(n=m^2\) vertices, \(|E|=2m(m-1)\) edges, and \(\Delta=4\) for interior vertices (\(\Delta=2\) or~\(3\) on the boundary). The chromatic number is~\(2\), since the grid is bipartite.

\item[Random \(d\)-regular graph.]
A graph chosen uniformly at random from all simple \(d\)-regular graphs on \(n\) vertices \cite{W99}. Every vertex has degree exactly \(d\), so \(\Delta=d\) and \(|E|=nd/2\). For \(d\geq3\), random regular graphs are typically well-connected with high girth relative to their size, making them a useful benchmark that avoids the structural regularity of lattice-based graphs.

\end{description}

\end{document}
